\documentclass[journal,onecolumn]{IEEEtran}

% correct bad hyphenation here
\hyphenation{op-tical net-works semi-conduc-tor re-inforce-ment simu-la-tion auto-ma-tion tech-niques}

\usepackage{hyperref}
\usepackage{xurl}

% put in your preamble
\usepackage{tabularx}

\usepackage{adjustbox}
\usepackage{amsmath}
\usepackage{amssymb}
\usepackage{caption}

\usepackage{booktabs}
\usepackage{multirow}
\usepackage{enumitem}

\captionsetup[table]{skip=4pt} % tighten caption spacing
\usepackage[numbers]{natbib}

% Widen the left and right margins a little (top/bottom left as IEEEtran default)
\usepackage[letterpaper, left=1.25in, right=1.25in]{geometry}

\begin{document}

% ===================================================================
% COVER PAGE
% ===================================================================
\begin{titlepage}
\centering

\vspace*{1cm}

{\Large\bfseries UNIVERSIDAD DE LOS ANDES}\\[0.3cm]
{\normalsize Faculty of Engineering}\\[0.1cm]
{\normalsize Department of Systems and Computing Engineering}\\

\vspace{2cm}

{\large\bfseries Flexible Policy Learning for Joint Activity\\[0.3cm]
Routing and Resource Allocation in\\[0.3cm]
Business Process Simulation}\\

\vspace{1.8cm}

{\normalsize\bfseries David Sequera}\\[0.2cm]
{\small d.sequera@uniandes.edu.co}\\

\vspace{1.5cm}

{\small A thesis submitted in partial fulfillment of the requirements\\[0.15cm]
for the degree of}\\[0.4cm]
{\normalsize\bfseries Magister en Ingeniería de Sistemas y Computación}\\

\vspace{1.8cm}

\begin{tabular}{@{}r@{\hspace{1em}}l@{}}
  \textbf{Advisor:}   & \'Oscar Fernando Gonz\'alez Rojas  \\[0.5cm]
  \textbf{Jury:}      & Manuel Alejandro Camargo Chavez    \\[0.2cm]
                      & Nicolas Cardozo Alvarez            \\
\end{tabular}

\vfill

{\small Bogot\'a D.C., 2026}

\end{titlepage}
% ===================================================================
% END COVER PAGE
% ===================================================================

 
\title{Flexible Policy Learning for Joint Activity Routing and Resource Allocation in Business Process Simulation}
 
\author{
  \IEEEauthorblockN{David Sequera}
  \IEEEauthorblockA{
    d.sequera@uniandes.edu.co\\
    Advisor: Oscar Fernando Gonz\'alez Rojas \\
    Universidad de los Andes, CSW
  }
}
 
\markboth{Flexible Policy Learning for Joint Activity Routing and Resource Allocation in Business Process Simulation
}%
{Sequera \MakeLowercase{\textit{et al.}}: Thesis Document}
 
\maketitle

\begin{abstract}
Business process optimization through simulation has advanced along two prominent
yet largely independent lines: resource allocation, which decides which resource executes
each activity, and prescriptive process monitoring, which decides when to make an intervention
in a case. Although both target overlapping operational objectives, their models and action spaces
remain separate, limiting the flexibility available to a single optimization approach.
This work presents a methodology that unifies both perspectives through a joint
decision framework based on deep reinforcement learning~(DRL). The proposed method
expands the action space by allowing the agent to jointly select, for a given case,
the next activity and its executing resource at each decision point, enabling it to
explore process configurations that neither approach can reach in isolation. We
formalize this problem as a Markov Decision Process over a stochastic simulation
environment learned from event logs, and define an architecture in which a DRL agent
learns a policy that maximizes compliance with service-level agreements~(SLAs),
defined as cycle time thresholds. To preserve realistic process behavior despite the
enlarged action space, the agent operates under a \textbf{flexible masked policy}
guided by learned branching probabilities, flexible enough to improve the process,
yet constrained enough to prevent greedy or unrealistic decisions. We introduce
\textbf{flexibility}, the capacity to dynamically delimit the action space in a way
that is tunable between process conformance and operational performance, as the
missing dimension in both lines of prior work. The approach is validated on three
event logs against heuristic and learning-based baselines, and evaluated
using both SLA compliance metrics and a multi-perspective similarity framework
covering control-flow, temporal, resource, and congestion dimensions. Results show
that jointly optimizing activity routing and resource allocation increases SLA
compliance while preserving a partially realistic process behavior across the 
measured perspectives.
\end{abstract}

\begin{IEEEkeywords}
Deep Reinforcement Learning, Business Process Simulation, Joint Optimization,
Activity Routing, Resource Allocation, Flexible Masked Policy,
Service-Level Agreements, Markov Decision Process.
\end{IEEEkeywords}


\IEEEpeerreviewmaketitle

\section{Introduction}
\label{sec:introduction}

The execution of a business process involves two intertwined decisions at every
step: which activity to perform next for a given case, and which resource should
perform it. In practice, these decisions are typically kept separate. Process
execution follows a control-flow model discovered from historical data, while
resource allocation is handled by a scheduling policy layered on top of that model.
This separation is operationally convenient; it allows the simulation to remain
faithful to the original control-flow, but it is strategically suboptimal. The
best resource assignment for a fixed routing may not achieve the performance of a
jointly selected activity-resource pair that routes the case differently from the
outset.

Recent research in both resource allocation and prescriptive process monitoring has
explored learning-based approaches to process optimization, though from largely
separate angles. On the resource allocation side, deep reinforcement learning~(DRL)
has emerged as a promising framework for learning assignment policies in business
process simulation~(BPS), with recent work demonstrating improvements over heuristic
baselines in cycle time reduction~\cite{Optimizing_Resource_Allocation_Policies,
Learning_policies_drl}. These approaches, however, uniformly fix the control-flow
and restrict the agent to selecting resources for activities determined by a separate
routing model. On the prescriptive process monitoring side, approaches such
as~\cite{ppm_optimal_policy_temporal_data, ppm_resource_constraints} explore
activity-level recommendations, but typically optimize case outcomes through
domain-specific interventions rather than through a general joint decision policy
over the full action space. Neither line of work addresses the joint selection of
activity and resource as a unified decision.

This work proposes a unified decision framework in which a DRL agent jointly
selects, at each decision point, the next activity and its executing resource as a
single action~$(a, r)$. The key challenge is that expanding the action space in this
way risks producing policies that achieve improvements by exploiting the simulator,
that is, by following paths that are fast in simulation but inconsistent with real
process behavior. The central contribution of this work is a \textbf{flexible
masked policy} that resolves this challenge: the mask grants the agent meaningful
flexibility to explore alternative process paths, while constraining its action space
to a process-consistent subset derived from learned branching probabilities. We
introduce \textbf{flexibility}, the capacity to dynamically delimit the action space
in a way that is tunable between process conformance and operational
performance, as the missing dimension in both lines of prior work.

The flexible masked policy is specified through three complementary filters: (i)~a
top-$k$ filter that limits the number of candidate activities; (ii)~a top-$p$ filter
that bounds the cumulative probability mass of selectable actions; and (iii)~a
minimum terminal probability threshold $p_{\mathrm{end}}^{\min}$, which prevents
the agent from closing a case prematurely before sufficient process progress has been
made. Together, these mechanisms preserve realistic routing behavior without
sacrificing the agent's capacity to explore alternative paths. We further introduce a
complementary regularization mechanism that shapes what the agent \emph{prefers} to
do within the admissible action set, drawing the learned policy toward the empirical
routing distribution without forbidding deviation. While the mask defines the
boundary of the action space, the regularization softens the agent's behavior within
that boundary, addressing a conformance challenge that is particularly acute in the
joint activity-resource setting where the enlarged action space makes the trade-off
between performance and realism more consequential than in prior work.

The contributions of this work are threefold. First, we formalize joint activity
routing and resource allocation as a Markov Decision Process~(MDP) over a stochastic
BPS environment learned from event logs, defining a case-aware state representation
that includes each case's positional prefix, SLA urgency, and branching probability
vector, alongside the global resource utilization features used in prior work. This
reframes the optimization problem and produces a unified decision framework where
prior work has treated routing and allocation as separate concerns. Second, we
introduce the flexible masked policy as the algorithmic mechanism that makes joint
optimization viable, formalizing flexibility as an explicit, configurable property
of the framework. Third, we introduce a regularization mechanism that incentivizes
the agent to prefer process-consistent actions within the admissible set, extending
the framework's ability to navigate the conformance-performance trade-off beyond what
the mask alone can achieve.

We validate the approach on three real-world event logs of varying scale and
domain, academic credential management and financial loan processing at two
scales, against five baselines forming a progressive ablation from fully random
decisions through heuristic and machine-learning-driven routing to DRL-based resource
allocation and, finally, to full joint DRL control. This design directly isolates the
contribution of each component of the proposed framework.

\subsection*{Research Objectives}

The general objective of this work is to design and validate a deep
reinforcement learning framework that jointly optimizes activity routing and
resource allocation in business process simulation, increasing compliance with
service-level agreements while preserving realistic process behavior. This
general objective is pursued through the following specific objectives:

\begin{itemize}
\item Formalize joint activity routing and resource allocation as a Markov
Decision Process over a stochastic business process simulation environment
learned from event logs, defining a case-aware state representation suitable for
the enlarged action space.
\item Design a flexible masked policy that constrains the agent's action space to
a process-consistent subset derived from learned branching probabilities, and
formalize flexibility as an explicit, configurable property tunable between
process conformance and operational performance.
\item Introduce a policy-level regularization mechanism that draws the learned
policy toward the empirical routing distribution without forbidding deviation,
and characterize its interaction with the masking mechanism.
\item Evaluate the framework on real-world event logs against a progressive
ablation of heuristic, machine-learning-driven, and DRL-based baselines, using
both SLA compliance metrics and a multi-perspective similarity framework covering
control-flow, temporal, resource, and congestion dimensions.
\end{itemize}

\subsection*{Summary of Validation Results}

The experimental validation provides consolidated evidence for the value of
joint optimization. The full joint policy (DRL-DRL) reaches near-perfect SLA
compliance at the training threshold $T_{75}$ across all three event logs
(CR\,$\geq$\,99.4\%), improving the compliance-improvement rate (CIR) by
$+32.8\%$ to $+33.3\%$ over the as-is process. Isolating the contribution of
joint control, the step from resource-only DRL (DM-DRL) to full joint DRL
roughly doubles the CIR gain (e.g., from $+20.4\%$ to $+32.8\%$ on
AC\_CRE and from $+11.7\%$ to $+33.3\%$ on BPIC17\_W); on BPIC12\_W, where every
Markov-based baseline performs below a random-routing policy, the joint agent
nonetheless recovers to $+33.2\%$ CIR. On the conformance side, the policy-level
regularization term lowers control-flow deviation (NGD) by up to $62\%$ relative
to the unregularized agent, at a compliance cost of at most $2.1$ percentage
points, confirming that the masking and regularization mechanisms are
complementary rather than competing.

The remainder of this article is structured as follows.
Section~\ref{sec:theoretical_framework} introduces the theoretical background
on business process simulation and reinforcement learning.
Section~\ref{sec:state_of_the_art} reviews related work and positions the
contribution relative to the two lines of prior research.
Section~\ref{sec:motivation_and_novelty} articulates the motivation and novelty
of the proposed approach.
Section~\ref{sec:research_questions} states the research questions and principal
research challenges.
Section~\ref{sec:problem_definition} formalizes the problem as a Markov Decision
Process and defines the notion of flexibility.
Section~\ref{sec:approach} presents the approach, covering the simulation
environment, state representation, flexible masked policy, regularization
mechanism, agent architecture, and reward function.
Section~\ref{sec:flex_evaluation} describes the evaluation design, datasets,
baselines, and metrics, and reports the results for the joint policy and the
flexibility mask.
Section~\ref{sec:reg_evaluation} evaluates the policy-level regularization term
through the mask\,$\times$\,KL ablation.
Section~\ref{sec:conclusion} concludes and outlines future
directions.

\section{Theoretical Framework}
\label{sec:theoretical_framework}

\subsection{Business Process Simulation}
\label{subsec:bps}

Business Process Simulation (BPS) is a practical approach to evaluating a
business process model before applying changes to the real process. Where
classical process analysis seeks closed-form analytical results, BPS creates
an \emph{executable digital twin} that mimics arrivals, work steps, queues,
and resource constraints to estimate cycle time, throughput, cost, and
service-level compliance under uncertainty. In contrast to purely analytical
queuing models, BPS embraces real-world complexities such as time-varying
demand, resource calendars, priorities, rework, and human behavior, making
it well-suited for what-if analysis and risk-aware decision-making
\cite{law_kelton_2015, banks_2010, laguna_marklund_2019, declarative}.

A process simulation model is typically represented as a flow using the
Business Process Model and Notation (BPMN). It is composed of activities
(tasks), resources (agents, machines, or teams), and gateways (transitions).
This flow is a directed graph of activities (nodes) and gateways (edges),
such as XOR and AND splits, that define the probability of transition from
one activity to another. A particular path through the graph is called a
\emph{trace}, and a single execution of an activity is called an \emph{event},
characterized by an associated activity, resource, start time, and end time.
To simulate a real process, the model draws on information such as resource
calendars, arrival time distributions, routing probabilities,
processing time and waiting time distributions and resource capacities.

One of the most established paradigms for BPS is \emph{discrete-event
simulation} (DES) \cite{law_kelton_2015, banks_2010, laguna_marklund_2019},
which reduces the simulation problem to its basic unit, the event, so that
the state of the process at any moment can be described as the set of events
currently being executed. When building a simulation model from historical
data, three broad approaches have emerged in the literature.

\emph{Data-Driven Process Simulation} (DDPS) defines the structural
components of the process (flow, resources, calendars, branching
probabilities, processing and waiting times) and discovers their
parameter values directly from an event log.
DDPS approaches make explicit modeling assumptions that aid
interpretability but may miss patterns that are difficult to express in the
chosen model structure. \emph{Deep learning-based simulation} replaces
hand-crafted components with machine learning models, capturing temporal and behavioral
patterns at the cost of reduced interpretability. \emph{Hybrid Process
Simulation} (HPS) combines the interpretable structure of DDPS with the
predictive accuracy of deep learning, typically using DDPS for the process
flow and neural models for time prediction \cite{review_bps,
camargo2023hybrid, Runtime_integration_of_simulation_for_business_processes}.

The present work adopts an empirical DDPS baseline as its simulation
foundation. Routing follows a first-order Markov policy in which branching
probabilities are estimated as conditional activity frequencies computed
directly from the event log. Processing times and case arrivals are sampled
empirically from the observed log distributions, preserving distributional
shape without parametric assumptions. Resource availability is governed by a
weekly calendar policy derived from observed resource activity patterns.
The DRL agent is layered on top of this empirical simulator, overriding the
routing and resource allocation interfaces at each decision point, while all
time-related components remain data-driven. The integration of ML-based time
and routing prediction, as realized in hybrid approaches such as
\cite{Runtime_integration_of_simulation_for_business_processes}, is
acknowledged as a valuable direction for future work.

\subsection{Reinforcement Learning}
\label{subsec:rl}

Reinforcement Learning (RL) is one of the three principal paradigms of
Machine Learning, alongside supervised and unsupervised learning. Unlike
supervised learning, which requires labeled input-output pairs, or
unsupervised learning, which discovers structure in unlabeled data, RL relies
on an agent that learns by interacting with an environment. At each step, the
agent observes a state, selects an action, and receives a reward signal that
indicates whether the action was beneficial. The agent's objective is to learn
a \emph{policy}, a mapping from states to actions, that maximizes the
cumulative expected reward over time \cite{sutton_2018_rl_intro}.

The mathematical foundation of RL is the \emph{Markov Decision Process}
(MDP), defined by a tuple $\mathcal{M} = (\mathcal{S}, \mathcal{A},
\mathcal{T}, \mathcal{R})$. The \textbf{state space} $\mathcal{S}$ describes
all configurations the agent may encounter. The \textbf{action space}
$\mathcal{A}$ defines the interactions available to the agent. The
\textbf{transition function} $\mathcal{T}(s' \mid s, a)$ encodes the
stochastic dynamics of the environment, specifying the probability of
reaching state $s'$ after taking action $a$ in state $s$. The \textbf{reward
function} $\mathcal{R}(s, a)$ quantifies the desirability of each transition,
encoding the operational objective. The MDP framework is appropriate when
outcomes are stochastic, depend on the current state, and cannot be evaluated
without simulating their downstream effects, precisely the conditions that
characterize business process execution.

RL algorithms can be broadly classified into value-based, policy-based, and
actor-critic approaches \cite{Optimizing_Resource_Allocation_Policies}. A
foundational value-based algorithm is \emph{Q-learning}, which estimates a
state-action value function $Q(s, a)$ representing the expected cumulative
reward of taking action $a$ in state $s$ and following the optimal policy
thereafter. Q values are updated iteratively using the Bellman equation,
which relates the value of a state to the expected values of its successor
states \cite{sutton_2018_rl_intro}. Q-learning is well-suited to discrete,
low-dimensional action spaces but scales poorly as the number of states and
actions grows.

A foundational policy-based algorithm is \emph{REINFORCE}, in which the
policy $\pi(a \mid s)$ is parameterized directly, typically by a neural
network. Rather than estimating values, REINFORCE learns from complete
trajectories of states, actions, and returns, updating the policy parameters
to increase the likelihood of actions that led to higher cumulative rewards
\cite{sutton_2018_rl_intro}. Policy-based methods handle continuous and
high-dimensional action spaces more naturally than value-based methods, but
suffer from high variance in gradient estimates and slow convergence.

\emph{Actor-Critic} methods synthesize the strengths of both approaches by
training two components jointly. The \textbf{actor} is a policy network that
selects actions; the \textbf{critic} is a value network that estimates the
expected future return from a given state. The critic's estimate is used to
compute the \emph{advantage}, defined as the difference between the actual
return and the critic's prediction, which indicates whether the actor's
decision was better or worse than expected. Learning from advantage rather
than raw returns reduces variance significantly and accelerates convergence
\cite{lapan_2024_drll_handson}. Both value-based and policy-based approaches
are viable for business process optimization; however, because actor-critic
architectures handle high-dimensional continuous state spaces efficiently and
underpin the most competitive modern DRL algorithms, this work adopts an
actor-critic design for its agent.

\subsection{Deep Reinforcement Learning and PPO}
\label{subsec:drl}

The emergence of deep neural networks marked a turning point for
reinforcement learning. Earlier approaches were limited by the need to
enumerate states and actions explicitly, restricting their applicability to
small, discrete domains. Deep neural networks, as function approximators,
allow both the policy and the value function to be represented over
arbitrarily large state and action spaces. This combination, known as
\emph{Deep Reinforcement Learning} (DRL), has produced remarkable results
in domains ranging from game playing to robotics, and has more recently been
applied to business process optimization
\cite{Optimizing_Resource_Allocation_Policies, Learning_policies_drl}.

Among DRL algorithms, \emph{Proximal Policy Optimization} (PPO)
\cite{lapan_2024_drll_handson} has emerged as one of the most widely adopted
actor-critic methods, combining sample efficiency, training stability, and
ease of implementation. PPO's central mechanism is a \emph{clipping
objective} that constrains the magnitude of each policy update, preventing
the new policy from deviating too far from the previous one. Formally, the
clipped surrogate objective is:
\[
\mathcal{L}^{\mathrm{CLIP}}(\theta)
=
\mathbb{E}_t\!\left[
\min\!\left(
r_t(\theta)\hat{A}_t,\;
\mathrm{clip}\!\left(
r_t(\theta),\,1-\varepsilon,\,1+\varepsilon
\right)\hat{A}_t
\right)
\right]
\]
where $r_t(\theta) = \pi_\theta(a_t \mid s_t) / \pi_{\theta_{\mathrm{old}}}
(a_t \mid s_t)$ is the probability ratio between the new and old policies,
$\hat{A}_t$ is the advantage estimate at time step $t$, and $\varepsilon$ is
the clipping hyperparameter. When the advantage is positive, the objective
rewards increasing the probability of the action, but only up to a factor of
$1 + \varepsilon$; when the advantage is negative, it penalizes the action,
but not beyond $1 - \varepsilon$. This prevents destructively large updates
that would collapse training, a particularly important property in the
long-horizon, non-stationary setting of business process simulation where
rewards are sparse at the episode level.

PPO estimates advantages using \emph{Generalized Advantage Estimation} (GAE),
which combines multi-step returns with exponential weighting to balance the
bias-variance trade-off in gradient estimation. GAE with parameters $\gamma$
(discount factor) and $\lambda$ (trace decay) produces smoother and more
reliable advantage estimates than single-step TD errors, improving the
quality of policy updates in environments with delayed and stochastic rewards.
In the business process setting, where the operational objective (SLA
compliance) is a case-level outcome that depends on a long sequence of
routing and allocation decisions, GAE plays a critical role in propagating
the terminal reward signal back through the trace to intermediate decision
points.

\subsection{Business Process Simulation as a Reinforcement Learning Problem}
\label{subsec:bps_as_rl}

The structural properties of business process simulation make it a natural
fit for the reinforcement learning framework. A business process unfolds as
a sequence of discrete decision points, at each of which a choice must be
made about which activity to execute next and which resource to assign to it.
The outcome of any such choice is stochastic: processing times, waiting
times, and downstream routing all depend on factors that cannot be fully
anticipated at the moment of decision. The operational objective, such as
minimizing cycle time or maximizing SLA compliance, is determined only at the
end of a complete case execution, creating a long-horizon credit assignment
problem that classical optimization methods struggle with.

These properties map directly onto the MDP abstraction. The state captures
the current configuration of the process, including resource utilization,
active cases, and case-specific features. The action encodes the decision to
be made at each step. The transition function reflects the stochastic
dynamics of the simulator. The reward encodes the operational objective,
providing the agent with a signal that connects individual decisions to
process-level outcomes. By framing BPS as an MDP, the problem becomes
amenable to DRL: the agent can learn, through repeated interaction with the
simulator, which routing and allocation decisions lead to better process
performance across the full distribution of cases and scenarios.

Prior work has exploited this connection primarily for resource allocation,
fixing the control-flow and training agents to assign resources within a
predetermined routing structure
\cite{Optimizing_Resource_Allocation_Policies, Learning_policies_drl,
Resource_allocation_using_task_similarity}. This work extends the connection
by expanding the action space to include both routing and resource allocation,
treating the joint decision as a single learnable policy. The theoretical
framework established in this chapter provides the foundation for that
extension; the following chapter surveys the existing literature to identify
precisely where the gap lies and why prior approaches cannot bridge it without the notion of flexibility introduced in this work.

\section{State of the Art}
\label{sec:state_of_the_art}

Business process optimization through simulation has developed along
two largely independent lines of work, each addressing a different
dimension of the joint decision problem this work seeks to unify.
This section surveys both lines, reviews the simulation environments
that underpin them, and closes with a characterization table that
positions the contribution relative to prior work.

\subsection{Simulation Environments for Process Optimization}
\label{subsec:soa_environments}

The quality of any learning-based process optimization approach
depends critically on the fidelity of the simulation environment
in which the agent is trained. If the environment does not represent
the real process accurately, the agent learns policies that exploit
simulation artifacts rather than genuine process structure. Two
components play distinct roles in this setting. The \textbf{simulator}
executes the business process faithfully taking into account calendar availability,
resources capabilities, advancing cases according
to a control-flow model and predicting processing and waiting times.
The \textbf{environment} wraps the simulator and mediates interaction
with the agent, translating the simulation state into an observable
representation, computing the reward signal, and enforcing the
action interface.

Data-driven approaches to building simulation environments have
converged on a core set of techniques for discovering these components
from event logs: a process flow representing the control structure,
resource calendars and skill assignments, empirical arrival and
service-time distributions, and branching probabilities estimated
from observed transition frequencies
\cite{law_kelton_2015, banks_2010, laguna_marklund_2019}. Tools
such as SIMOD \cite{SIMOD} automate the discovery of these components
using statistical and process mining techniques, producing simulation
models ready for direct use as training environments.

A limitation of purely data-driven simulators is that their branching
probabilities are typically estimated as first-order conditional
frequencies, which may not capture the dependence of routing decisions
on the full history of a case.
\citeauthor{Runtime_integration_of_simulation_for_business_processes}~\cite{Runtime_integration_of_simulation_for_business_processes}
address this through the RIMS framework, which integrates machine
learning predictions of processing and waiting times at runtime
during simulation rather than as a post-processing step. RIMS
proposes, among its artifacts, decision trees for next-activity
prediction and deep learning models for time estimation, enabling
queue-aware predictions that improve simulation accuracy in congested
processes.
\citeauthor{camargo2023hybrid}~\cite{camargo2023hybrid} propose a
complementary hybrid approach that uses process mining to discover
the control-flow structure and deep learning to predict processing
and waiting times, retaining partial what-if analysis capabilities
while improving temporal accuracy. Both works represent state-of-the-art
approaches in simulation fidelity and suggest that hybrid architectures
can offer meaningful improvements over purely data-driven baselines.

The quality of simulation models has been studied systematically by
\citeauthor{framework_for_measuring_the_quality_of_business_process_simulation_models}~\cite{framework_for_measuring_the_quality_of_business_process_simulation_models},
who propose a multi-perspective evaluation framework comparing
simulated logs against a reference log across control-flow, temporal,
resource, and congestion dimensions. This framework provides the
evaluation methodology adopted in the present work and establishes
reference values against which the behavioral realism of DRL-generated
logs can be assessed.

\subsection{Resource Allocation Optimization}
\label{subsec:soa_resource_allocation}

The dominant approach in the resource allocation line treats the
routing policy as fixed and restricts the optimization agent to
assigning resources to activities as they become ready for execution.
Reinforcement learning has been applied to this problem with
encouraging results across several process structures.

\citeauthor{Learning_policies_drl}~\cite{Learning_policies_drl}
propose two complementary methods for resource allocation in business
process simulation: a DRL policy that learns assignment decisions
directly, and a score-based value function approximation that learns
feature weights for prioritizing resource assignments. Both methods
are evaluated on six archetypal process structures and three composite
processes, outperforming or matching heuristic baselines in the
majority of scenarios. The routing policy in this work is determined
by the simulation model and the agent controls only the resource
assignment interface.

\citeauthor{Optimizing_Resource_Allocation_Policies}~\cite{Optimizing_Resource_Allocation_Policies}
combine DRL with a hybrid simulation model to evaluate resource
allocation on real-world processes, reporting improvements over
heuristic baselines in 14 out of 24 evaluation scenarios, with gains
reaching up to 307\% in specific configurations. Their work
demonstrates that DRL-based resource allocation generalizes beyond
synthetic processes; the routing policy, however, remains outside
the scope of the agent's control.

\citeauthor{Resource_allocation_using_task_similarity}~\cite{Resource_allocation_using_task_similarity}
approach resource allocation from a different angle, applying
reinforcement learning to reduce worklist entropy by exploiting
task similarity. Their premise is that assigning similar tasks to
the same resource reduces cycle time by amortizing shared subtasks.
An entropy measure over worklists serves as the optimization
objective, and RL searches for resource assignments that minimize
it. In this approach, the activity sequence is also not subject to
the agent's control.

\citeauthor{genetic_resource_allocations}~\cite{genetic_resource_allocations}
take a non-RL approach, applying genetic algorithms to discover
optimal resource configurations for what-if scenarios by simulating
candidate configurations and evaluating their cost, flow time,
waiting time, and workload trade-offs. While this approach considers
multi-objective optimization and what-if scenarios, the routing
policy remains determined by the simulation model throughout.

Across all of these approaches, the action space is uniformly limited
to the resource policy interface. The routing of cases through the
process is determined by the simulation model and cannot be intervened
upon by the agent.

\subsection{Prescriptive Process Monitoring}
\label{subsec:soa_prpm}

Prescriptive Process Monitoring (PrPM) approaches operate at the
activity level but typically frame decisions as domain-specific
interventions rather than as general routing choices over the full
action space.

\citeauthor{ppm_resource_constraints}~\cite{ppm_resource_constraints}
propose a PrPM method that learns intervention policies under
resource constraints using DRL combined with conformal prediction.
The agent decides when to trigger a predefined domain-specific
intervention for a given case, optimizing case outcomes rather than
cycle time. The action space is binary: intervene or not. The work
demonstrates that incorporating uncertainty estimates from conformal
prediction into the reward signal accelerates convergence toward
higher-gain intervention policies, but the framework is designed
for a specific intervention type and does not generalize to arbitrary
routing or allocation decisions.

\citeauthor{ppm_optimal_policy_temporal_data}~\cite{ppm_optimal_policy_temporal_data}
construct an MDP over trace clusters and apply policy iteration to
recommend optimal next actions for ongoing cases. Cases are clustered
using prefix representations based on activity frequencies and
positions, and the MDP is solved with Monte Carlo simulation to
estimate state values. The action space is defined over abstract
state-cluster pairs rather than concrete activity-resource pairs,
which limits the granularity of the recommendations and prevents
direct control over activities assignment, but provides the agent
a broader action space allowing not a specific intervention but a map
of actions that can lead to a better KPI optimization.

\subsection{Additional Perspectives}
\label{subsec:soa_additional}

Beyond the two main lines of work, several contributions provide
useful context for understanding the broader landscape of business
process optimization.

\citeauthor{Optimization_of_Activity_Batching_Policies}~\cite{Optimization_of_Activity_Batching_Policies}
address a related but distinct problem: optimizing the grouping and
activation of activity instances through batching. Their Pareto
optimization framework balances waiting time, processing effort, and
cost using DRL combined with intervention heuristics, and represents
an example of how RL can be applied to process-level structural
decisions beyond resource assignment.

\citeauthor{The_Business_Process_Optimization_Competition}~\cite{The_Business_Process_Optimization_Competition}
present a process optimization competition in which participants
develop algorithms to decide the timing and sequencing of patient
admissions in a hospital process with uncertain treatment tasks.
Eight algorithms are evaluated, including simulated annealing,
evolutionary algorithms, tabu search, and DRL. Although this work
is tailored to a specific medical domain and does not address the
general joint activity-resource decision problem, it provides an
insightful comparison of different optimization methods, highlighting
RL as a particularly effective paradigm for time-aware scenarios.

\citeauthor{Leveraging_Counterfactuals}~\cite{Leveraging_Counterfactuals}
propose a prescriptive framework that leverages counterfactual generation
to recommend joint activity--resource pairs for running process instances.
By employing DiCE to efficiently explore alternative next-step combinations,
the framework avoids the suboptimality of sequential approaches that fix
the activity first and assign a resource independently. A configurable
reduction threshold further allows analysts to balance the proportion
of cases receiving recommendations against the magnitude of the improvements
achieved. This work shares the joint decision framing but differs in
two respects: the policy is learned over a full stochastic simulation
where outcomes reflect the interaction of all concurrent cases rather
than single-case predictions, and routing deviation from observed process
behavior is explicitly bounded through a tunable masking mechanism.


From the perspective of emerging techniques,
\citeauthor{PM_LLM_Benchmark}~\cite{PM_LLM_Benchmark} evaluate the
ability of large language models to perform process mining tasks,
including process description and insight extraction.
\citeauthor{business_process_modeling_llms}~\cite{business_process_modeling_llms}
evaluate the generation of process models from natural language
descriptions. Both works suggest that embedding representations can
capture business process semantics in ways that may be relevant to
future state representations for DRL agents, and are noted here as
a promising direction rather than a direct predecessor to this work.

\subsection{Characterization of Related Work}
\label{subsec:soa_characterization}

Table~\ref{tab:CharacterizationTable} summarizes the most closely
related approaches along seven dimensions. \textbf{Technique} refers
to the learning or optimization method used. \textbf{Optimization
Metric} is the operational measure being optimized. \textbf{Testing}
describes how results are compared against a reference. \textbf{Action
Space} describes what the approach controls at each decision point;
this is the most diagnostic dimension and the one that most clearly
reveals the gap. \textbf{Simulation Type} indicates whether the
approach operates on the process as observed (AS IS) or considers
hypothetical modifications (What-if). \textbf{Trace-Aware} indicates
whether the agent's state includes information about the history of
the current case instance; trace-aware approaches can condition their
decisions on the prefix of the ongoing trace, while those that are
not treat each decision point independently of prior events in the
same case. \textbf{Application Field} classifies the contribution
as belonging to resource allocation, prescriptive process monitoring,
or a related area. The optimization scope column is omitted from the
table as all surveyed approaches operate at the whole-process level.

The table reveals a consistent pattern across both lines of work.
Resource allocation approaches uniformly restrict the action space
to the resource policy interface; the routing of cases is determined
by the simulation model and cannot be intervened upon by the agent.
PrPM approaches reach the activity level but frame decisions as
domain-specific interventions or abstract cluster transitions rather
than as general activity-resource pairs. No prior approach jointly
controls both the next activity and its executing resource as a
unified, general-purpose decision. This is the precise gap that the
present work addresses.

\begin{table*}[!htbp]
\centering
\footnotesize
\setlength{\tabcolsep}{2pt}
\renewcommand{\arraystretch}{1.05}
\begin{tabularx}{\textwidth}{|>{\raggedright\arraybackslash}p{2.0cm}
                             |>{\raggedright\arraybackslash}p{1.3cm}
                             |>{\raggedright\arraybackslash}p{1.2cm}
                             |>{\raggedright\arraybackslash}p{1.3cm}
                             |>{\raggedright\arraybackslash}p{1.3cm}
                             |>{\raggedright\arraybackslash}p{1.2cm}
                             |>{\raggedright\arraybackslash}p{0.9cm}
                             |>{\raggedright\arraybackslash}p{1.2cm}
                             |>{\raggedright\arraybackslash}X|}
\hline
\textbf{Article} &
\textbf{Tech.} &
\textbf{Opt. Metric} &
\textbf{Testing} &
\textbf{Action Space} &
\textbf{Sim.} &
\textbf{Trace-Aw.} &
\textbf{Field} &
\textbf{Purpose} \\ \hline
\textbf{\citeauthor{Optimizing_Resource_Allocation_Policies}~\cite{Optimizing_Resource_Allocation_Policies}} &
DRL & Cycle Time & Heuristics & Resources & What-if & Yes & Resource Alloc. &
Optimize resource allocation via activity--resource pairs \\ \hline
\textbf{\citeauthor{Learning_policies_drl}~\cite{Learning_policies_drl}} &
DRL, Bayes. Opt. & Cycle Time & Heuristics & Resources & AS IS & No & Resource Alloc. &
Learn adaptive resource allocation policies across process structures \\ \hline
\textbf{\citeauthor{Resource_allocation_using_task_similarity}~\cite{Resource_allocation_using_task_similarity}} &
RL & Cycle Time & Heur.\ \& Log & Resources & AS IS & No & Resource Alloc. &
Reduce worklist entropy via task similarity \\ \hline
\textbf{\citeauthor{genetic_resource_allocations}~\cite{genetic_resource_allocations}} &
Genetic Alg. & Multiobj. & What-if Scen. & Resources & What-if & No & Resource Alloc. &
Discover optimal resource configurations via GA \\ \hline
\textbf{\citeauthor{Optimization_of_Activity_Batching_Policies}~\cite{Optimization_of_Activity_Batching_Policies}} &
DRL & Cycle Time \& Cost & Heuristics & Activity Batch & AS IS & No & Activity Batch. &
Optimize batching of activity instances via Pareto meta-heuristics \\ \hline
\textbf{\citeauthor{ppm_resource_constraints}~\cite{ppm_resource_constraints}} &
DRL & Case Outcome & DRL Base. & Binary Interv. & AS IS & No & PrPM &
Learn uncertainty-aware intervention policies under resource constraints \\ \hline
\textbf{\citeauthor{ppm_optimal_policy_temporal_data}~\cite{ppm_optimal_policy_temporal_data}} &
RL & Generic KPI & DRL Base. & Trace Clust., Action & AS IS & No & PrPM &
Recommend optimal next actions via MDP policy iteration over trace clusters \\ \hline
\textbf{\citeauthor{Leveraging_Counterfactuals}~\cite{Leveraging_Counterfactuals}} &
Grad. Boost. & Cycle Time & Seq. Base. & Activities \& Res. & No, trace Reconstruction & Yes & PrPM &
Recommend joint activity--resource pairs via counterfactual search over a cycle-time predictor \\ \hline
\textbf{This work} &
DRL (PPO) & Cycle Time (SLA) & Heur.\ \& Log & Activities \& Res. & AS IS \& What-if & Yes & RA \& PrPM &
Learn a flexible joint policy over activity routing and resource allocation via a tunable masked policy \\ \hline
\end{tabularx}
\caption{Characterization of related work. All surveyed approaches
operate at the whole-process level. The Action Space column is the
most diagnostic dimension: resource allocation approaches restrict
control to the resource policy interface, while PrPM approaches
operate on interventions or abstract cluster transitions. To the best of our knowledge, no prior work in process optimization has jointly controlled activity routing and resource allocation through a tunable masked policy learned over a full stochastic simulation.}
\label{tab:CharacterizationTable}
\end{table*}

\subsection{Positioning and Research Gap}
\label{subsec:soa_gap}

The two lines of work reviewed above share a common limitation,
though it manifests differently in each. Resource allocation
approaches offer no flexibility over routing: the control-flow
cannot be intervened upon by the RL agent, and decisions are
confined to the resource policy interface
\cite{Learning_policies_drl, Optimizing_Resource_Allocation_Policies,
Resource_allocation_using_task_similarity, genetic_resource_allocations}.
PrPM approaches do reason about which activity should occur next,
but without a principled mechanism to bound how far the recommended
path may deviate from observed process behavior
\cite{ppm_resource_constraints, ppm_optimal_policy_temporal_data}.
In both cases, the trade-off between process realism and optimization
potential is not tunable: it is either fully resolved in favor of
the original routing policy, or left largely unconstrained.

We argue that \emph{flexibility}, understood as the capacity to
dynamically delimit the action space in a way that is tunable
between process conformance and operational improvement, is the
missing dimension in both lines of work. A decision framework that
incorporates flexibility as a first-class property can interpolate
between conservative behavior that closely mirrors the original
process and more exploratory behavior that pursues performance gains
through alternative routing, with the degree of exploration governed
by interpretable parameters. To the best of our knowledge, no prior
work in business process optimization has formalized this notion
of flexibility as part of the decision framework, nor proposed a
mechanism to implement it within a learned policy. The following
chapters address that gap.

\section{Motivation and Novelty}
\label{sec:motivation_and_novelty}

Business process optimization through simulation has long operated
under a structural assumption that is rarely questioned: that routing
and resource allocation are separate concerns, each addressed by a
different model. Resource allocation approaches treat the control-flow
as given and restrict the agent to the resource policy interface,
while prescriptive process monitoring approaches reason about which
activity should occur next through domain-specific interventions that
are not designed to interact with the resource layer. The hypothesis underlying
joint optimization is that simultaneously controlling routing and allocation
exposes configurations hardly reachable by either decision independently.

The central motivation of this work is that jointly controlling
activity routing and resource allocation can unlock process
configurations that improve operational performance beyond what
either approach achieves independently. A resource allocation agent
cannot reroute a case toward a shorter path; its decisions are
confined to assigning resources to activities whose sequence has
already been determined. A prescriptive monitoring approach that
recommends alternative activities cannot simultaneously optimize
which resource executes them. The performance ceiling of each
approach is therefore set not by the quality of its learning
algorithm but by the boundaries of its action space. Expanding
those boundaries is a structural change that requires a proper framing.

The missing ingredient is a formal notion of \emph{flexibility}:
the capacity to dynamically delimit the action space in a way that
is tunable between process conformance and operational performance.
Flexibility is a desirable property of any framework that operates
over such a broad action space, serving as a hyper-parametrization
of the recommendations between realism and optimization. Without it,
an unconstrained agent will tend to exploit the simulator, discovering
paths that improve simulation metrics while becoming progressively
detached from real process behavior. Resource allocation researchers
never needed to bound routing exploration because routing was never
part of their action space. Prescriptive monitoring researchers,
operating through predefined intervention types, did not require a
general mechanism to constrain arbitrary routing choices. The joint
frontier is where the problem first becomes visible, and where
flexibility becomes necessary.

This work introduces flexibility as a first-class property of
the decision framework, defined through a \textbf{flexible
masked policy} that grounds the boundary of the admissible action
space in learned branching probabilities. The mask exposes three
interpretable control parameters: a top-$k$ filter on the number
of candidate activities, a top-$p$ filter on the cumulative
probability mass of selectable actions, and a minimum terminal
probability threshold $p_{\mathrm{end}}^{\min}$ that prevents
premature case closure. The novelty lies not in any individual
component but in the formalization of flexibility as a configurable,
interpretable property that transforms the conformance-performance
trade-off from an implicit modeling assumption into an explicit
design choice.

The work further introduces a complementary regularization
mechanism that shapes the agent's preferences within the admissible
action set, drawing the learned policy toward the empirical routing
distribution without forbidding deviation. Where the mask defines
what the agent may do, the regularization shapes what the agent
tends to do, addressing the conformance challenge at a finer
granularity. Together, they form a coherent system for navigating
the conformance-performance trade-off across both the boundary
and the interior of the action space.


\section{Research Questions}
\label{sec:research_questions}


The work seeks to address the following research questions:

\begin{quote}
\textit{\textbf{RQ1.} Can a deep reinforcement learning agent,
operating on a stochastic business process simulation, learn a
decision policy for activity routing and resource allocation that
increases the proportion of cases completing within a predefined
service-level agreement (SLA) threshold?}

\textit{\textbf{RQ2.} How does the behavior induced by the learned
DRL policy differ from realistic process executions across
control-flow, temporal, congestion, and resource dimensions?}

\textit{\textbf{RQ3.} What is the magnitude of the improvement
achieved by the full DRL-based decision policy compared to purely
heuristic and purely machine-learning-driven baselines, and how
does performance degrade when the agent's scope of control is
restricted?}

\textit{\textbf{RQ4.} To what extent can a policy-level
regularization term reduce the control-flow differences introduced
by the DRL routing policy without degrading SLA compliance?}
\end{quote}

RQ1 targets the primary operational objective: whether a jointly
trained policy can genuinely improve SLA compliance over a stochastic
simulation environment. RQ2 addresses the realism dimension,
examining whether the agent's decisions produce logs that remain
behaviorally consistent with the original process across multiple
perspectives. RQ3 isolates the contribution of each component of
the framework through a progressive ablation, connecting the
experimental design directly to the contributions of the work.
RQ4 examines the regularization term: whether drawing the policy
toward the as-is routing distribution narrows the control-flow gap
the routing agent introduces while preserving its SLA gains.

\section{Problem Definition}
\label{sec:problem_definition}

This chapter formalizes the optimization problem addressed by this
work. It introduces the formal primitives of a business process,
defines the decision point as the unit of interaction between the
agent and the simulation, maps those primitives onto the components
of a Markov Decision Process, and closes with a summary statement
that positions the problem in its operational context.

\subsection{Process Model and Formal Primitives}
\label{subsec:formal_primitives}

We begin by establishing the formal vocabulary used throughout the
remainder of this work.

\begin{itemize}

    \item \textbf{Activity Set} (\(A\)):
    An activity set is a finite set of operations that can be
    performed within a process:
    \[
        A = \{ a_1, a_2, \dots, a_m \}.
    \]

    \item \textbf{Resource} (\(r \in R\)):
    A resource is an entity qualified to execute a subset of
    activities:
    \[
        r \in R, \quad A_r \subseteq A,
    \]
    where $R$ is the finite set of all resources and $A_r$ is the
    set of activities that resource $r$ is qualified to perform.

    \item \textbf{Event} (\(e\)):
    An event is an activity instance represented as the tuple
    \[
        e = (c,\, a,\, t_s,\, t_e,\, r),
    \]
    where $c \in C$ is the case identifier, $a \in A$ is the
    executed activity, $t_s, t_e \in \mathbb{R}_{\geq 0}$ are the
    start and end timestamps with $t_s \leq t_e$, and $r \in R$
    is the executing resource. The set of all events is denoted
    $E$.

    \item \textbf{Trace} (\(\sigma\)):
     A trace is the unique sequence of activities that characterizes
    a particular execution path through the process:
    \[
        \sigma = \langle a_1, a_2, \dots, a_\ell \rangle,
        \quad a_i \in A.
    \]
    Multiple cases may share the same trace if they follow the
    same sequence of activities. The prefix of a trace up to
    position $k$ is denoted $\sigma_{:k} = \langle a_1, \dots,
    a_k \rangle$.

    \item \textbf{Event Log} (\(L\)):
    An event log is a finite collection of events grouped by
    case, from which the set of observed traces and their
    empirical frequencies can be derived:
    \[
        L = \{ e_1, e_2, \dots, e_n \}, \quad e_i \in E.
    \]

    \item \textbf{Case} (\(c \in C\)):
    A case is a single instance of a process execution. Multiple
    cases may follow the same trace; each case is independently
    associated with a set of performance metrics, including
    processing time, waiting time, resource usage, and cycle
    time, where cycle time is defined as
    $t_e(e_{|\sigma_c|}) - t_s(e_1)$.

    \item \textbf{Decision Point}:
    A decision point arises at any branching point in the process
    where the current execution state of a case $c$ admits more
    than one valid continuation. At each decision point, two
    interdependent choices must be made: which activity $a \in A$
    to execute next, and which resource $r \in R$ to assign to it.
    Prior work resolves these choices separately, fixing activity
    selection through a routing model and delegating only resource
    assignment to the optimization policy. This work treats them
    as a single joint action $\alpha = (a, r)$.

\end{itemize}

\subsection{MDP Formulation}
\label{subsec:mdp_formulation}

We formalize the optimization problem as a Markov Decision Process
(MDP) $\mathcal{M} = (\mathcal{S}, \mathcal{A}, \mathcal{T},
\mathcal{R})$. This framework is appropriate here because the
outcomes of routing and resource allocation decisions are stochastic,
depend on the current process state, and cannot be evaluated without
simulating their downstream effects on the full case lifecycle.

The \textbf{state space} $\mathcal{S}$ describes all configurations
the agent may encounter during a simulation run. Each state $s \in
\mathcal{S}$ captures a fixed-size snapshot of the process at a
given decision point: which cases are active, how long each has been
running, what activity each is executing or waiting for, and which
resources are busy or available, including a case-specific prefix
that summarizes the routing history of the case under consideration.

The \textbf{action space} $\mathcal{A}$ defines the joint decisions
available to the agent. An action $\alpha = (a, r) \in \mathcal{A}$
is the joint selection of a next activity and an executing resource
for a given case at a given decision point, constrained by resource
qualification so that only pairs $(a, r)$ with $a \in A_r$ are
admissible.

The \textbf{transition function} $\mathcal{T}(s' \mid s, \alpha)$
encodes the stochastic dynamics of the simulation environment.
Given a state $s$ and a joint action $\alpha$, it specifies the
probability distribution over next states $s'$. In the BPS setting,
this reflects how routing and allocation decisions propagate through
the simulation: a particular choice of activity and resource may
reduce waiting time, free a bottleneck resource, or accelerate case
completion, but the exact outcome depends on the state of the
broader process, including the processing times and arrival patterns
of other active cases.

The \textbf{reward function} $\mathcal{R}(s, \alpha)$ encodes the
operational objective. The agent receives two types of reward
signal: a soft, dense signal at each decision step that provides
directional feedback on the current cycle time relative to the SLA
threshold, and a hard terminal signal issued at the moment a case
completes that encodes the compliance outcome definitively. The
formal specification of both components is given in
Section~\ref{subsec:reward}.

\subsection{Problem Statement}
\label{subsec:problem_statement}

Given an event log $L$ from a business process, the goal is to
learn a joint decision policy $\pi: \mathcal{S} \rightarrow
\mathcal{A}$ that, at each decision point of each active case,
selects a joint action $\alpha = (a, r)$ so as to maximize the
proportion of cases whose cycle time falls below a
practitioner-defined threshold $T$. The inputs to the problem
are the event log $L$; a simulation model derived from it that
encodes the process structure, resource pool, time distributions,
and empirical routing probabilities across six policy interfaces
(described in Section~\ref{subsec:simulation_environment}); and
the threshold $T$. The decision variable is the joint action
$\alpha = (a, r)$ at each decision point, jointly selecting the
next activity to execute and the resource to assign to it. The
objective is to maximize SLA compliance, formalized through the
reward function specified in Section~\ref{subsec:reward}.

The threshold $T$ is implemented in this work as a
percentile of the empirical cycle time distribution of $L$,
producing a synthetic but well-grounded target. The framework
is not specific to this choice: any objective expressible as
a reward signal, whether based on cost, waiting time, resource
utilization, or a composite KPI, can be substituted without
modifying the agent or the environment architecture.


\section{Approach}
\label{sec:approach}

This chapter presents the framework, describing the simulation
environment, state representation, flexible masked policy, agent
architecture, reward function, and regularization mechanism that
together implement the MDP formulation of
Section~\ref{sec:problem_definition}. Figure~\ref{img:C4Context}
provides a system-level view of the five containers that compose
this system: the \textbf{Initializer}, which bootstraps the simulation
model from the event log; the \textbf{Discrete-Event Simulator},
which executes the process faithfully according to its policy
interfaces; the \textbf{Simulation Environment}, which mediates
the interaction between the simulator and the agent; the
\textbf{RL Agent}, which learns the joint activity-resource
policy; and the \textbf{Evaluator}, which assesses the quality
of the generated logs. The process analyst interacts with the
system by configuring and launching the Initializer and by
analyzing the results produced by the Evaluator.

\begin{figure}[!ht]
  \centerline{\includegraphics[width=\columnwidth]{images/c4_context.png}}
  \caption{C4 context view of the system. The process analyst
  interacts with the Initializer to configure the simulation and
  with the Evaluator to analyze results. The five containers,
  Initializer, Discrete-Event Simulator, Simulation Environment,
  RL Agent, and Evaluator, form the complete optimization pipeline.}
  \label{img:C4Context}
\end{figure}

\subsection{Simulation Environment}
\label{subsec:simulation_environment}

The architecture distinguishes two components with complementary
responsibilities, illustrated in Figure~\ref{img:Architecture}.
The \textbf{Discrete-Event Simulator} executes the business process
faithfully: it advances cases through the process structure,
samples processing and waiting times from data-driven distributions,
and maintains the state of all active cases and resources. The
\textbf{Simulation Environment} wraps the simulator and mediates
the interaction with the DRL agent. It is responsible for four
functions: computing the flexible mask $\mathcal{F}_c$ at each
decision point; translating the simulation state into a vector
observable by the agent; translating the agent's output into
concrete activity and resource assignments that the simulator can
execute; and computing the reward signal from the resulting case
outcomes. This separation keeps the simulator self-contained and
reusable: the environment layer can be modified, for example to
change the reward function or state representation, without
altering the underlying simulation logic.

\begin{figure}[!ht]
  \centerline{\includegraphics[width=\columnwidth]{images/env_architecture.png}}
  \caption{Component-level architecture of the Simulation Environment
  and RL Agent. The environment exposes four functions to the agent:
  state representation, action masking, action selection, and reward
  computation. The agent's two output heads, activity and resource,
  feed back into the masking layer before a joint action is executed
  in the simulator.}
  \label{img:Architecture}
\end{figure}

The simulator is built around a policy interface architecture that
decouples each behavioral dimension of the process from the
simulation engine itself. Six interfaces define the simulator's
full behavioral surface.

\begin{itemize}
    \item \textbf{Routing Policy}: determines the next activity
    at each decision point, implemented as a first-order Markov
    policy in which branching probabilities are estimated as
    conditional activity frequencies from the event log.

    \item \textbf{Processing Time Policy}: determines the
    duration of each activity execution, sampled empirically
    per activity from the observed log distributions, preserving
    distributional shape without parametric assumptions.

    \item \textbf{Waiting Time Policy}: models delays before
    execution begins, implemented via extraneous waiting times
    stratified by activity-resource pair.

    \item \textbf{Arrival Policy}: controls how new cases enter
    the process, sampled empirically from the observed
    inter-arrival distribution.

    \item \textbf{Calendar Policy}: governs resource availability
    over time using a weekly grid derived from observed resource
    activity patterns, determining availability via a
    percentile-based threshold per weekday-hour slot.

    \item \textbf{Resource Policy}: determines which resource is
    assigned to a given activity based on skill compatibility.
\end{itemize}

Any implementation conforming to a given interface can be
substituted without modifying the simulation engine, making the
architecture extensible by design. At each decision point, the
routing and resource interfaces are overridden by the DRL agent
through the environment layer: rather than sampling the next
activity from the branching distribution, the environment exposes
the decision to the agent and uses the branching probabilities to
compute $\mathcal{F}_c$ as described in
Section~\ref{subsec:mask}. All time-related interfaces remain
data-driven throughout training and evaluation.

\subsection{State Representation}
\label{subsec:state_representation}

The state $s \in \mathbb{R}^{d}$ is a fixed-size vector observed
by the agent at each decision point, normalized to $[0, 1]$ to
improve sample efficiency~\cite{sutton_2018_rl_intro}. It
concatenates four groups of features that together capture the
global process configuration, the demand for each activity, the
case-specific context, and the temporal environment.

Let $R = \{r_1, \dots, r_{|R|}\}$ be the ordered set of
resources, $A = \{a_1, \dots, a_{|A|}\}$ the ordered set of
activities, $\sigma_c = \langle e_1, \dots, e_n \rangle$ the
prefix of case $c$, and $T$ the SLA threshold. The constants
$\theta_q$ and $\theta_\lambda$ are normalization caps for
queue-based features and trace length respectively, calibrated
empirically across the event logs used in the evaluation.

\begin{itemize}
    \item \textbf{Utilization} ($u_i$): the fraction of resource
    $r_i$'s capacity currently in use:
    \[
        u_i = \frac{\text{count}(r_i)}{\text{cap}(r_i)}.
    \]

    \item \textbf{Assignment} ($\eta_i$): an encoding of the
    activity that resource $r_i$ is currently executing:
    \[
        \eta_i =
        \begin{cases}
            j / |A|, & \text{if } r_i \text{ is executing } a_j, \\
            0,       & \text{if } r_i \text{ is idle.}
        \end{cases}
    \]

    \item \textbf{Queue Pressure} ($q_i$): the number of cases
    waiting for resource $r_i$, normalized and clamped:
    \[
        q_i = \min\!\left(\frac{|\text{queue}(r_i)|}{\theta_q},\;
        1\right).
    \]

    \item \textbf{Pending Demand} ($\kappa_j$): the number of
    cases awaiting execution of activity $a_j$:
    \[
        \kappa_j = \min\!\left(
        \frac{|\text{pending}(a_j)|}{\theta_q},\; 1\right).
    \]

    \item \textbf{Last Activity} ($\ell_c$): the position of
    case $c$ in the process:
    \[
        \ell_c =
        \begin{cases}
            j / |A|, & \text{if the last executed activity is }
            a_j, \\
            0, & \text{if } \sigma_c = \langle\rangle.
        \end{cases}
    \]

    \item \textbf{Trace Length} ($\lambda_c$): the number of
    activities executed so far in case $c$:
    \[
        \lambda_c = \min\!\left(\frac{|\sigma_c|}{\theta_\lambda},
        \; 1\right).
    \]

    \item \textbf{SLA Urgency} ($\phi_c$): the elapsed time of
    case $c$ relative to the SLA threshold:
    \[
        \phi_c = \min\!\left(\frac{t_{\text{now}} -
        t_s(e_1)}{T},\; 1\right).
    \]

    \item \textbf{Branching Probabilities} ($\mathbf{b}_c$): the
    transition probability vector from the routing policy at the
    current position of case $c$:
    \[
        \mathbf{b}_c = \big[\, P(a_j \mid a_{\text{last}})
        \,\big]_{j=1}^{|A|}.
    \]
    These components play a dual role: they inform the agent's
    routing decisions and are used directly to compute
    $\mathcal{F}_c$ in the masked policy.

    \item \textbf{Hour of Day} ($\tau_h$) and \textbf{Day of
    Week} ($\tau_d$): normalized temporal features that allow
    the agent to learn calendar-dependent
    patterns~\cite{Optimizing_Resource_Allocation_Policies}:
    \[
        \tau_h = \frac{t_{\text{now}} \bmod 86400}{86400},
        \qquad
        \tau_d = \frac{\lfloor t_{\text{now}} / 86400 \rfloor
        \bmod 7}{6}.
    \]
\end{itemize}

The full state vector concatenates all components:
\begin{align*}
    s = \big[\,
    &\underbrace{u_1, \eta_1, q_1,\; \dots,\;
    u_{|R|}, \eta_{|R|}, q_{|R|}}_{\text{global: } 3|R|}, \\
    &\underbrace{\kappa_1, \dots,
    \kappa_{|A|}}_{\text{demand: } |A|},\;\;
    \underbrace{\ell_c, \lambda_c, \phi_c,\;
    \mathbf{b}_c}_{\text{case: } |A|+3},\;\;
    \underbrace{\tau_h, \tau_d}_{\text{time: } 2}
    \,\big],
\end{align*}
with dimension $d = 3|R| + 2|A| + 5$. Resource availability is
not enforced as a hard constraint at action selection time;
rather, the utilization and queue pressure features provide the
agent with sufficient information to learn availability-aware
assignment through the reward signal, keeping the action space
definition clean and the policy general.



\subsection{The Flexible Masked Policy}
\label{subsec:mask}

The flexible masked policy is the central algorithmic contribution
of this work and the mechanism through which flexibility, as
defined in Section~\ref{sec:motivation_and_novelty}, is
implemented. At each decision point, before the agent selects an
action, the set of candidate activities is restricted to a
process-consistent subset $\mathcal{F}_c \subseteq A$ derived
from the branching probability vector $\mathbf{b}_c$. Three
complementary filters define this subset conjunctively: an
activity must satisfy all applicable conditions simultaneously
to be admitted into $\mathcal{F}_c$.

\begin{itemize}
    \item \textbf{Top-$k$ filter}: retains only the $k$
    activities with the highest branching probability, limiting
    routing choices to the most process-consistent options.

    \item \textbf{Top-$p$ filter}: retains the minimal set of
    activities whose cumulative probability mass reaches or
    exceeds $p$, ensuring that retained candidates collectively
    account for a sufficient proportion of observed routing
    behavior.

    \item \textbf{Minimum terminal probability threshold}
    ($p_{\mathrm{end}}^{\min}$): applies exclusively to terminal
    activities, enforcing a lower bound on the probability of
    any selectable terminal activity. This prevents the agent
    from closing a case prematurely before sufficient process
    progress has been made. Let $\Omega \subseteq A$ denote the
    set of terminal activities.
\end{itemize}

Formally, the admissible activity set is:
\[
\mathcal{F}_c =
\left\{
a_j \in A \;\middle|\;
\begin{aligned}
& a_j \in \mathrm{top}\text{-}k(\mathbf{b}_c), \\
& a_j \in \mathrm{top}\text{-}p(\mathbf{b}_c), \\
& a_j \notin \Omega \;\lor\;
  P(a_j \mid a_{\mathrm{last}}) \ge p_{\mathrm{end}}^{\min}
\end{aligned}
\right\}.
\]

Actions outside $\mathcal{F}_c$ receive a logit value of
$-10^9$ prior to the softmax (a numerical stand-in for $-\infty$), following standard masking
practice in DRL~\cite{Learning_policies_drl}. The parameters
$k$, $p$, and $p_{\mathrm{end}}^{\min}$ together control the
degree of flexibility: lower values produce a conservative policy
that closely mirrors the original routing distribution; higher
values expand the agent's freedom to explore alternative paths
at the cost of greater divergence from observed behavior. For
the resource dimension, a complementary mask enforces skill
compatibility, restricting the selectable resource set to
$R_{\text{qual}}(a) = \{r \in R : a \in A_r\}$ for the chosen
activity $a$.

\subsection{Agent Architecture}

The agent is an Actor-Critic network trained with Proximal Policy
Optimization (PPO). Its defining feature is a hierarchical, two-head
action selection design: the actor first selects an activity
$a \in \mathcal{F}_c$ from the masked activity distribution, then
selects a resource $r \in R_{\text{qual}}(a)$ from a resource
distribution conditioned on the chosen activity via an embedding.
This conditioning is what makes the joint action $\alpha = (a, r)$
genuinely joint: the resource head sees the selected activity as
context, so the two components are chosen together rather than
independently. This hierarchical coupling of routing and allocation
within a single policy is, to our knowledge, the principal architectural
departure of this work from prior agents that learn only one of the two
decisions. The critic takes the same state $s$ as input and outputs a
scalar estimate $V(s)$ of the expected future return, used as a baseline
to compute the advantage of each action. Both actor and critic share a
common trunk of fully connected layers with separate output heads, and
both the activity and resource masks are applied as a $-10^9$ logit fill
before sampling, the standard masking approach in
DRL~\cite{Learning_policies_drl}.

The PPO clipping objective constrains the magnitude of each policy update,
preventing any single update from moving the policy too far from the previous
one. This stability is what makes PPO well suited to the long-horizon, dynamic
control that this environment demands, where many cases evolve concurrently
and each decision plays out over a long sequence of downstream steps, but it
also places the burden on the objective: a stable learner optimizes faithfully
whatever it is told to optimize, and little more. With compliance encoded
through the intermediate and terminal rewards, the agent reliably reaches high
SLA performance; what it does not do is express any preference among the many
paths that achieve that performance equally well, because nothing in the
objective asks it to. In an environment where strong performance is reachable
along behaviorally different paths, defining the objective correctly,
including what to preserve and not only what to improve, is therefore as
consequential as the choice of learning algorithm. This observation motivates
the regularized objective developed in Section~\ref{sec:reg_evaluation}.


\subsection{Reward Function}
\label{subsec:reward}

The reward function guides the agent toward cycle time compliance through a
piecewise-linear design with two regimes: an \textbf{intermediate shaping
signal} issued at every step, and a \textbf{terminal gate} that encodes SLA
compliance at case completion. Let $ct(\sigma)$ denote the elapsed cycle time of case $\sigma$,
$T$ the service-level threshold, and $H$ a scaling constant controlling the magnitude of rewards
and penalties, set to $H=1$ in all experiments.

During execution, the shaping signal provides dense directional feedback:

\[
\mathcal{R}_{\text{int}}(s,\alpha)=
\begin{cases}
\dfrac{H}{1000}\!\left(1-\dfrac{ct(\sigma)}{T}\right), &
  \text{if } ct(\sigma)<T,\\[10pt]
-\dfrac{H}{100}\cdot\dfrac{ct(\sigma)}{T}, &
  \text{if } ct(\sigma)\ge T.
\end{cases}
\]

When the elapsed time remains within the SLA budget, the agent receives a
positive reward that decreases linearly from $+H/1000$ (at case start) to $0$
as $ct(\sigma)\to T$, continuously favoring faster progress. Once the threshold
is exceeded, the signal turns negative and grows in magnitude proportionally to
the degree of violation, increasing pressure to avoid further delay. At case
completion, the terminal reward overrides the shaping signal:

\[
\mathcal{R}_{\text{term}}(s,\alpha)=
\begin{cases}
+H, & ct(\sigma)<T,\\
-H, & ct(\sigma)\ge T.
\end{cases}
\]

The $H/1000$ scaling factor on the intermediate signal is deliberate: it ensures
that the terminal reward dominates the learning objective, so the agent
optimizes for SLA compliance as an outcome rather than promoting the agent
to stay in the simulation for a bigger reward.
The intermediate component acts as a shaping term in the sense of
reward shaping~\cite{sutton_2018_rl_intro}, providing denser gradient
information without distorting the compliance objective.

\subsection{Regularized Training Objective}
\label{subsec:regularization_method}

The approximation developed so far equips the agent with several
mechanisms to reach the desired outcome: a well-defined objective, a
reward function that encodes SLA compliance, and a flexible mask that
restricts the action space to plausible activities. None of these,
however, conveys to the agent that staying close to the original
process behavior is itself a goal. The agent is free to satisfy the
SLA through any admissible path, including paths that meet the
threshold while drifting away from the routing patterns observed in
the as-is log. To couple performance with behavioral similarity, we
introduce this pressure not as an additional reward term but as a
term in the training objective itself. The agent is trained to
minimize
%
\[
\begin{aligned}
\mathcal{L}(\theta) =\;
& \underbrace{-\,\mathcal{L}^{\mathrm{CLIP}}(\theta)
  + c_v\,\mathcal{L}^{\mathrm{VF}}(\theta)
  - c_e\,\mathcal{H}[\pi_\theta]}_{\text{SLA optimization (PPO)}} \\[4pt]
& +\;
  \lambda_{\text{KL}}\,
  \underbrace{\mathbb{E}_s\!\left[
  \mathrm{KL}\!\left(
  \tilde\pi_\theta(\cdot\mid s)\,\big\|\,
  \tilde P(\cdot\mid a_{\text{last}},\mathcal{F}_c)
  \right)\right]}_{\text{similarity pressure}}
\end{aligned}
\]
%
where the first group is the standard PPO loss that drives SLA
compliance, and the second is the forward KL divergence between the
agent's masked activity distribution $\tilde\pi_\theta$ and the
empirical routing distribution $\tilde P$, both restricted and
renormalized over the feasible set $\mathcal{F}_c$. The latter acts as
a regularization term: it draws the policy toward the as-is routing
distribution, which serves as the reference for the process behavior
we aim to preserve. The coefficient $\lambda_{\text{KL}}\ge 0$ sets the
strength of this term; $\lambda_{\text{KL}}=0$ recovers the
unregularized objective.

The two groups pull in tension at every parameter update: the PPO term
pushes the policy toward whichever activity improves SLA compliance,
while the KL term pulls it back toward the as-is routing distribution.
Because the KL gradient is applied directly to the policy, the agent
can move away from the routing prior only where the PPO gradient is
locally strong enough to overcome it; that is, only where deviating
measurably improves the KPI. The result is \emph{selective deviation}:
the agent conforms wherever the competing routing alternatives lead to
similar outcomes, and spends its deviation budget only on the changes
that are necessary to meet the SLA. The objective does not suppress
deviation; it forces the agent to justify each deviation by a
performance gain.

This is the property the flexibility mask alone cannot provide. The
mask restricts the easy outs, removing implausible activities from
$\mathcal{F}_c$, but it does not present the agent with the need to
identify the decision points where an intervention is genuinely
required to change the outcome of the KPI. Within the admissible set,
an unregularized agent is free to favor the least similar option
whenever it is marginally faster, with nothing to penalize the choice.
Filtering bounds what the agent \emph{may} do; it does not
\emph{incentivize} what the agent should prefer. The KL term supplies
this missing incentive, shaping the agent's preferences inside the
boundary the mask defines. 

\section{Flexibility Evaluation}
\label{sec:flex_evaluation}


This section evaluates RQ1--RQ3 defined
in Section~\ref{sec:research_questions}; RQ4 is addressed separately in
Section~\ref{sec:reg_evaluation}. RQ1 examines whether
the joint policy improves SLA compliance; RQ2 assesses the
behavioral realism of the learned policy across four process perspectives;
and RQ3 isolates the contribution of each component through a progressive
ablation against heuristic and learning-based baselines.


\subsection{Experimental Setup}
\label{subsec:setup}

\paragraph{Datasets.}
We evaluate on three real-world event logs
(Table~\ref{tab:datasets}): an academic credentials management
process (AC\_CRE) with high resource variability and low case
volume; a loan application subprocess from the BPIC 2012 challenge
(BPIC12\_W); and a loan application subprocess from the BPIC 2017
challenge (BPIC17\_W); both BPIC logs retaining only work-item
activities with start and complete timestamps; all logs
pre-processed by
\citeauthor{framework_for_measuring_the_quality_of_business_process_simulation_models}~\cite{framework_for_measuring_the_quality_of_business_process_simulation_models}.

\begin{table}[!ht]
\centering\scriptsize
\setlength{\tabcolsep}{4pt}
\begin{tabular}{lrrrrr}
\toprule
\textbf{Log} & \textbf{Cases} & \textbf{Events} &
\textbf{Activities} & \textbf{Resources} & \textbf{Traces} \\
\midrule
AC\_CRE   & 398     & 1{,}945  & 16 & 306 &     54 \\
BPIC12\_W & 3{,}030 & 16{,}338 &  6 &  47 &    735 \\
BPIC17\_W & 7{,}402 & 53{,}332 &  7 & 105 & 1{,}843 \\
\bottomrule
\end{tabular}
\caption{Training event logs (BPIC12\_W and BPIC17\_W retain only
W-activities with start and complete timestamps).}
\label{tab:datasets}
\end{table}

\paragraph{Policies.}
Every policy is an \emph{(ActivitySelector, ResourceSelector)} pair
executed for each log (Table~\ref{tab:policies}).
RA-RR is a random lower bound; DM-RR samples the empirical
first-order Markov routing; DM-GR pairs DM routing with a greedy
minimum-processing-time heuristic; DM-DRL pairs DM routing with a
PPO-trained resource-only agent; and DRL-DRL is the full joint
agent. Both DRL variants were trained at threshold $T_{75}$ and
evaluated at all four thresholds to assess generalization.

\begin{table}[!ht]
\centering\scriptsize
\setlength{\tabcolsep}{3pt}
\begin{tabular}{llll}
\toprule
\textbf{Policy} & \textbf{Activity} & \textbf{Resource} &
\textbf{Role} \\
\midrule
RA-RR   & Random uniform   & Random             & Lower bound \\
DM-RR   & Empirical Markov & Random             & As-is routing \\
DM-GR   & Empirical Markov & Min.\ proc.\ time  & Heuristic baseline \\
DM-DRL  & Empirical Markov & PPO (res.\ only)   & Resource-only DRL \\
DRL-DRL & PPO (joint)      & PPO (joint)        & Full joint DRL \\
\bottomrule
\end{tabular}
\caption{Evaluation policies. All share the same feasibility masks.}
\label{tab:policies}
\end{table}

\paragraph{Training configuration and metrics.}
Table~\ref{tab:training} reports the flexibility parameters and SLA
training threshold per log. All three logs share threshold $T_{75}$
and cumulative-probability filter $p=0.9$; the per-log choices of
$k$ and $p_{\mathrm{end}}^{\min}$ reflect the routing structure and
case-completion patterns of each process. For all metrics we run
$K=10$ simulated logs and report mean $\pm$ 95\%~CI.

\begin{table}[!ht]
\centering\small
\setlength{\tabcolsep}{5pt}
\begin{tabular}{lcccc}
\toprule
\textbf{Log} & \textbf{Threshold} & $k$ & $p$ &
$p_{\mathrm{end}}^{\min}$ \\
\midrule
AC\_CRE   & $T_{75}$ & 2 & 0.9 & 0.3 \\
BPIC12\_W & $T_{75}$ & 2 & 0.9 & 0.3 \\
BPIC17\_W & $T_{75}$ & 3 & 0.9 & 0.1 \\
\bottomrule
\end{tabular}
\caption{Flexibility parameters per log.}
\label{tab:training}
\end{table}

We evaluate operational performance via the compliance rate (CR),
the fraction of cases completing within threshold $T$, and the
compliance improvement rate (CIR), the relative CR gain over the
original log. Behavioral similarity is measured across four
perspectives~\cite{framework_for_measuring_the_quality_of_business_process_simulation_models}:
control-flow (NGD, $n$-gram distance on activity sequences);
temporal (AED, absolute hourly event distribution; RED, event
distribution relative to case start; CED, circadian pattern per
weekday); resource (CWD, circadian workforce distribution per
hour); and congestion (CAR, hourly case arrivals; CTD, cycle-time
histogram distance). Table~\ref{tab:metrics} details their
implementation.

\begin{table}[!ht]
\centering\small
\setlength{\tabcolsep}{3pt}
\begin{tabular}{@{} l l p{4.4cm} @{}}
\toprule
\textbf{Perspective} & \textbf{Metric} & \textbf{Method} \\
\midrule
\multirow{2}{*}{Operational}
  & CR  & $|\{\sigma:ct(\sigma)<T\}|/|\mathcal{L}|$ \\
  & CIR & $(\mathrm{CR}^{\mathrm{sim}}-\mathrm{CR}^{\mathrm{ref}})/\mathrm{CR}^{\mathrm{ref}}$ \\
\midrule
Control-flow & NGD & $n$-gram distance; SAE \\
\midrule
\multirow{3}{*}{Temporal}
  & AED & EMD on hourly event time-series \\
  & CED & Avg.\ EMD per weekday (circadian) \\
  & RED & EMD on events relative to case start \\
\midrule
Resource & CWD & Avg.\ EMD of active resource counts \\
\midrule
\multirow{2}{*}{Congestion}
  & CAR & EMD on hourly case arrivals \\
  & CTD & 1-Wasserstein on cycle-time histogram \\
\bottomrule
\end{tabular}
\caption{Evaluation metrics~\cite{framework_for_measuring_the_quality_of_business_process_simulation_models}.
Similarity metrics: lower is better ($\downarrow$).}
\label{tab:metrics}
\end{table}


\subsection{Performance Analysis}
\label{subsec:performance}

Table~\ref{tab:compliance_t75} reports CR and CIR at the training
threshold $T_{75}$ for all five policies, with policies as rows and
logs as columns so that the ablation progression reads top to
bottom within each log.

% ------------------------------------------------------------------
%  TABLE: CR/CIR at T75 — transposed (policies as rows, logs as cols)
% ------------------------------------------------------------------
\begin{table}[!ht]
\centering\scriptsize
\setlength{\tabcolsep}{2.5pt}
\renewcommand{\arraystretch}{1.1}
\begin{tabular}{ll ccc}
\toprule
\textbf{Policy} & \textbf{Met.} &
\textbf{AC\_CRE} & \textbf{BPIC12\_W} & \textbf{BPIC17\_W} \\
\midrule
\multirow{2}{*}{RA-RR}
  & CR  & 66.3{\tiny(±1.7)} & 79.0{\tiny(±0.3)} & 91.9{\tiny(±0.3)} \\
  & CIR & $-$11.4{\tiny(±2.3)} & +5.3{\tiny(±0.4)} & +22.6{\tiny(±0.4)} \\
\midrule
\multirow{2}{*}{DM-RR}
  & CR  & 71.0{\tiny(±1.3)} & 65.0{\tiny(±0.6)} & 80.1{\tiny(±0.4)} \\
  & CIR & $-$5.1{\tiny(±1.7)} & $-$13.3{\tiny(±0.8)} & +6.8{\tiny(±0.5)} \\
\midrule
\multirow{2}{*}{DM-GR}
  & CR  & 63.5{\tiny(±1.8)} & 61.6{\tiny(±0.6)} & 62.4{\tiny(±0.5)} \\
  & CIR & $-$15.2{\tiny(±2.4)} & $-$17.8{\tiny(±0.8)} & $-$16.8{\tiny(±0.7)} \\
\midrule
\multirow{2}{*}{DM-DRL}
  & CR  & 90.2{\tiny(±1.2)} & 65.6{\tiny(±0.4)} & 83.8{\tiny(±0.3)} \\
  & CIR & +20.4{\tiny(±1.6)} & $-$12.5{\tiny(±0.5)} & +11.7{\tiny(±0.4)} \\
\midrule
\multirow{2}{*}{DRL-DRL}
  & CR  & \underline{99.4{\tiny(±0.2)}} & \underline{99.9{\tiny(±0.0)}} & \underline{100.0{\tiny(±0.0)}} \\
  & CIR & \underline{+32.8{\tiny(±0.3)}} & \underline{+33.2{\tiny(±0.1)}} & \underline{+33.3{\tiny(±0.0)}} \\
\bottomrule
\end{tabular}
\caption{CR (\%) and CIR (\%) at $T_{75}$ (training target).
Values: mean (±95\% CI). Higher is better ($\uparrow$).
Best per (log, metric) pair underlined.}
\label{tab:compliance_t75}
\end{table}

\noindent\textbf{SLA compliance (RQ1).}
Table~\ref{tab:compliance_all} extends the view to all four
thresholds for DM-DRL and DRL-DRL. DRL-DRL achieves near-perfect
compliance at $T_{75}$ across all three logs
(CR\,$\geq$\,99.4\%; CIR\,+32.8\% to +33.3\%).
CIR reaches +5.2--5.3\% at $T_{95}$ and up to +99.2\% at $T_{50}$
on BPIC17\_W. This breadth is a consequence of the reward design:
the intermediate shaping signal provides a dense gradient for
cycle-time reduction as a continuous objective, so the agent learns
to accelerate cases broadly rather than concentrating them near the
training threshold. Resource-only DRL (DM-DRL) yields meaningful
gains where the Markov routing is adequate; up to +20.4\% CIR at
$T_{75}$ on AC\_CRE, but cannot saturate the SLA goal on its own,
confirming that routing control is necessary for full compliance.

% ------------------------------------------------------------------
%  TABLE: CR/CIR all thresholds — transposed
%  Rows: threshold × metric  |  Columns: log × policy
% ------------------------------------------------------------------
\begin{table*}[!ht]
\centering\scriptsize
\setlength{\tabcolsep}{3pt}
\renewcommand{\arraystretch}{1.1}
\begin{tabular}{ll cc cc cc}
\toprule
& &
\multicolumn{2}{c}{\textbf{AC\_CRE}} &
\multicolumn{2}{c}{\textbf{BPIC12\_W}} &
\multicolumn{2}{c}{\textbf{BPIC17\_W}} \\
\cmidrule(lr){3-4}\cmidrule(lr){5-6}\cmidrule(lr){7-8}
\textbf{Thr.} & \textbf{Met.} &
\textbf{DM-DRL} & \textbf{DRL-DRL} &
\textbf{DM-DRL} & \textbf{DRL-DRL} &
\textbf{DM-DRL} & \textbf{DRL-DRL} \\
\midrule
\multirow{2}{*}{$T_{95}$}
  & CR  & 98.4{\tiny(±0.4)} & \underline{99.9{\tiny(±0.1)}} & 84.8{\tiny(±0.5)} & \underline{100.0{\tiny(±0.0)}} & 97.4{\tiny(±0.2)} & \underline{100.0{\tiny(±0.0)}} \\
  & CIR & +3.6{\tiny(±0.4)} & \underline{+5.2{\tiny(±0.1)}} & $-$10.7{\tiny(±0.6)} & \underline{+5.3{\tiny(±0.0)}} & +2.5{\tiny(±0.2)} & \underline{+5.3{\tiny(±0.0)}} \\
\midrule
\multirow{2}{*}{$T_{90}$}
  & CR  & 97.1{\tiny(±0.6)} & \underline{99.7{\tiny(±0.1)}} & 79.0{\tiny(±0.7)} & \underline{100.0{\tiny(±0.0)}} & 94.2{\tiny(±0.3)} & \underline{100.0{\tiny(±0.0)}} \\
  & CIR & +8.0{\tiny(±0.7)} & \underline{+10.8{\tiny(±0.2)}} & $-$12.3{\tiny(±0.8)} & \underline{+11.1{\tiny(±0.0)}} & +4.7{\tiny(±0.3)} & \underline{+11.1{\tiny(±0.0)}} \\
\midrule
\multirow{2}{*}{$T_{75}^{\,\dagger}$}
  & CR  & 90.2{\tiny(±1.2)} & \underline{99.4{\tiny(±0.2)}} & 65.6{\tiny(±0.4)} & \underline{99.9{\tiny(±0.0)}} & 83.8{\tiny(±0.3)} & \underline{100.0{\tiny(±0.0)}} \\
  & CIR & +20.4{\tiny(±1.6)} & \underline{+32.8{\tiny(±0.3)}} & $-$12.5{\tiny(±0.5)} & \underline{+33.2{\tiny(±0.1)}} & +11.7{\tiny(±0.4)} & \underline{+33.3{\tiny(±0.0)}} \\
\midrule
\multirow{2}{*}{$T_{50}$}
  & CR  & 66.0{\tiny(±1.3)} & \underline{94.0{\tiny(±0.9)}} & 35.3{\tiny(±0.7)} & \underline{97.2{\tiny(±0.1)}} & 62.1{\tiny(±0.5)} & \underline{99.6{\tiny(±0.0)}} \\
  & CIR & +32.0{\tiny(±2.7)} & \underline{+88.0{\tiny(±1.8)}} & $-$29.3{\tiny(±1.4)} & \underline{+94.4{\tiny(±0.3)}} & +24.3{\tiny(±1.0)} & \underline{+99.2{\tiny(±0.1)}} \\
\bottomrule
\end{tabular}
\caption{CR (\%) and CIR (\%) across all thresholds for DM-DRL and
DRL-DRL. Values: mean (±95\% CI). Higher is better ($\uparrow$).
Best per row underlined.
$^\dagger$$T_{75}$ is the training target; remaining thresholds
test generalization.}
\label{tab:compliance_all}
\end{table*}

\noindent\textbf{Incremental benefit of joint control (RQ3).}
The five-policy ablation isolates each component's contribution.
On AC\_CRE and BPIC17\_W, the step from DM-DRL (+20.4\%, +11.7\%)
to DRL-DRL (+32.8\%, +33.3\%) raises the CIR by +12.4~pp and
+21.6~pp respectively; since both agents share the same simulator,
the difference is attributable entirely to joint routing control.

\noindent\textbf{BPIC12\_W anomaly.}
On BPIC12\_W, DM-RR, DM-GR, and DM-DRL all produce negative CIR
($-$13.3\%, $-$17.8\%, $-$12.5\%), meaning even random routing
(RA-RR, +5.3\%) outperforms every Markov-based policy.
This reflects a calibration limitation: BPIC12\_W likely contains
parallel structures that the XOR-only routing model approximates
with bidirectional transition probabilities, introducing artificial
loops that systematically favor slower paths.
DRL-DRL is unaffected; it learns its own routing from reward
feedback, and recovers to +33.2\% CIR, making this the strongest
single demonstration of the value of joint learning in this
evaluation.

\noindent\textbf{DM-GR as a negative control.}
The greedy processing-time heuristic consistently underperforms
DM-RR ($-$16.8\% CIR on BPIC17\_W), confirming that routing all
work to the fastest individual resources creates congestion that
erodes end-to-end times and that congestion-aware learned policies
are necessary.

% ---- 6.3  SIMILARITY ANALYSIS -------------------------------
\subsection{Similarity Analysis}
\label{subsec:similarity}

Table~\ref{tab:similarity_ref} anchors simulator fidelity against two reference frames. The upper block positions DM-RR among established simulators: DM-RR (train) reports the distance to the log from which the model was built; DM-RR (test) reports it against the held-out split, which was never observed during simulator construction; it is included here solely for comparison. On CTD, DM-RR falls below both SIMOD and ServiceMiner on AC\_CRE and BPIC17\_W, indicating closer cycle-time reproduction; NGD is higher across all three logs, a structural consequence of the XOR-only routing approximation, which cannot reproduce every control-flow path present in the data.
The lower block provides a natural scale anchor. Train vs Test reports the distance between the two temporal splits of the same event log, capturing the intrinsic variability of the process itself. Simulator distances of a similar order of magnitude suggest the model reproduces the process at a fidelity comparable to the drift between data splits. As simulator validation is not the focus of this work, this comparison should be interpreted as an indicative reference rather than a strict benchmark.

\begin{table}[!ht]
\centering\small
\setlength{\tabcolsep}{3pt}
\begin{tabular}{llcc}
\toprule
\textbf{Log} & \textbf{Source} & \textbf{NGD} & \textbf{CTD} \\
\midrule
\multirow{5}{*}{AC\_CRE}
& DM-RR (train) & 0.26{\tiny(±0.00)} & 34.2{\tiny(±4.1)} \\
& DM-RR (test)  & 0.24{\tiny(±0.01)} & 50.6{\tiny(±4.6)} \\
& SIMOD         & 0.24{\tiny(±0.01)} & 62.2{\tiny(±1.7)} \\
& ServiceMiner  & 0.13{\tiny(±0.01)} & 99.9{\tiny(±0.3)} \\
\cmidrule(lr){2-4}
& Train vs Test & 0.17               & 26.2              \\
\midrule
\multirow{5}{*}{BPIC12\_W}
& DM-RR (train) & 0.36{\tiny(±0.00)} & 170.3{\tiny(±9.6)} \\
& DM-RR (test)  & 0.34{\tiny(±0.00)} & 142.3{\tiny(±9.6)} \\
& SIMOD         & 0.56{\tiny(±0.00)} & 93.5{\tiny(±0.7)}  \\
& ServiceMiner  & 0.13{\tiny(±0.01)} & 124.2{\tiny(±0.3)} \\
\cmidrule(lr){2-4}
& Train vs Test & 0.08               & 28.0               \\
\midrule
\multirow{5}{*}{BPIC17\_W}
& DM-RR (train) & 0.18{\tiny(±0.00)} & 44.9{\tiny(±1.8)} \\
& DM-RR (test)  & 0.16{\tiny(±0.00)} & 15.2{\tiny(±1.8)} \\
& SIMOD         & 0.37{\tiny(±0.00)} & 102.9{\tiny(±0.8)} \\
& ServiceMiner  & 0.06{\tiny(±0.00)} & 112.8{\tiny(±0.2)} \\
\cmidrule(lr){2-4}
& Train vs Test & 0.08               & 30.9               \\
\bottomrule
\end{tabular}
\caption{DM-RR vs.\ SIMOD and ServiceMiner
\cite{framework_for_measuring_the_quality_of_business_process_simulation_models}
on NGD and CTD ($\downarrow$ better). Lower block shows reference
distances between data splits.}
\label{tab:similarity_ref}
\end{table}

Table~\ref{tab:similarity_drl} compares similarity metrics across
all evaluated policies, including \textit{+no flex}: DRL-DRL with
the activity mask removed.

\noindent\textbf{Behavioral similarity (RQ2).}
On AC\_CRE and BPIC17\_W, the mask reduces NGD by nearly half
relative to the unmasked agent (0.46 vs.\ 0.89 and 0.68 vs.\
0.84) at no compliance cost, demonstrating that flexibility
effectively trades off behavioral similarity against performance
when the routing prior accurately represents the process. On
BPIC12\_W this effect is absent, as the XOR-only approximation
of parallel structures distorts the branching probabilities the
mask relies on. The effectiveness of the flexibility mechanism is
therefore contingent on the fidelity of the underlying process
representation.

% ------------------------------------------------------------------
%  TABLE: Similarity — rows=log×policy, cols=metrics
%  Kept as table* for width; grouped by log with \midrule
% ------------------------------------------------------------------
\begin{table*}[!ht]
\centering\scriptsize
\setlength{\tabcolsep}{3pt}
\renewcommand{\arraystretch}{1.1}
\begin{tabular}{ll ccccccc}
\toprule
& &
\multicolumn{7}{c}{\textbf{Metric}} \\
\cmidrule(lr){3-3}\cmidrule(lr){4-5}\cmidrule(lr){6-7}\cmidrule(lr){8-9}
\textbf{Log} & \textbf{Policy}
  & \textbf{NGD}
  & \textbf{AED} & \textbf{RED}
  & \textbf{CED} & \textbf{CWD}
  & \textbf{CAR} & \textbf{CTD} \\
\midrule
\multirow{4}{*}{AC\_CRE}
  & DM-RR
    & \underline{0.26}$^{\pm.00}$
    & 76.7$^{\pm13.9}$ & \underline{22.8}$^{\pm1.6}$
    & \underline{2.03}$^{\pm.10}$ & \underline{2.07}$^{\pm.11}$
    & \underline{89.8}$^{\pm15.7}$ & \underline{34.2}$^{\pm4.1}$ \\
  & DM-DRL
    & 0.27$^{\pm.00}$
    & \underline{64.7}$^{\pm10.5}$ & 28.9$^{\pm2.8}$
    & 2.53$^{\pm.14}$ & 2.47$^{\pm.16}$
    & 90.0$^{\pm14.2}$ & 86.3$^{\pm5.8}$ \\
  & DRL-DRL
    & 0.46$^{\pm.00}$
    & 70.1$^{\pm16.2}$ & 70.4$^{\pm.7}$
    & 2.71$^{\pm.19}$ & 2.68$^{\pm.20}$
    & 88.3$^{\pm15.4}$ & 165.8$^{\pm1.6}$ \\
  & ~~+no flex
    & 0.89$^{\pm.01}$
    & 65.6$^{\pm11.1}$ & 71.3$^{\pm1.8}$
    & 2.33$^{\pm.22}$ & 2.35$^{\pm.20}$
    & 105.6$^{\pm18.1}$ & 176.2$^{\pm1.5}$ \\
\midrule
\multirow{4}{*}{BPIC12\_W}
  & DM-RR
    & \underline{0.36}$^{\pm.00}$
    & \underline{137.7}$^{\pm15.1}$ & \underline{118.8}$^{\pm6.9}$
    & \underline{2.18}$^{\pm.07}$ & \underline{1.44}$^{\pm.04}$
    & \underline{169.4}$^{\pm23.2}$ & 170.3$^{\pm9.6}$ \\
  & DM-DRL
    & 0.36$^{\pm.00}$
    & 123.0$^{\pm13.2}$ & 103.8$^{\pm5.2}$
    & 2.10$^{\pm.06}$ & 1.37$^{\pm.03}$
    & 170.8$^{\pm21.9}$ & 154.4$^{\pm7.2}$ \\
  & DRL-DRL
    & 0.75$^{\pm.00}$
    & 253.6$^{\pm23.8}$ & 165.0$^{\pm.3}$
    & 2.70$^{\pm.08}$ & 2.18$^{\pm.13}$
    & 176.2$^{\pm23.3}$ & \underline{146.5}$^{\pm.1}$$^{\dagger}$ \\
  & ~~+no flex
    & \underline{0.64}$^{\pm.00}$
    & 252.0$^{\pm23.4}$ & 166.4$^{\pm.4}$
    & 2.76$^{\pm.10}$ & 2.17$^{\pm.11}$
    & 170.9$^{\pm23.3}$ & 145.9$^{\pm.4}$ \\
\midrule
\multirow{4}{*}{BPIC17\_W}
  & DM-RR
    & \underline{0.18}$^{\pm.00}$
    & \underline{87.8}$^{\pm22.8}$ & \underline{11.5}$^{\pm2.0}$
    & \underline{1.53}$^{\pm.02}$ & \underline{0.78}$^{\pm.02}$
    & \underline{131.6}$^{\pm34.7}$ & \underline{44.9}$^{\pm1.8}$ \\
  & DM-DRL
    & 0.18$^{\pm.00}$
    & 96.7$^{\pm17.5}$ & 29.1$^{\pm2.2}$
    & 1.47$^{\pm.02}$ & 0.76$^{\pm.02}$
    & 131.4$^{\pm33.1}$ & 71.9$^{\pm3.4}$ \\
  & DRL-DRL
    & 0.68$^{\pm.00}$
    & 229.8$^{\pm44.1}$ & 179.1$^{\pm.2}$
    & 1.67$^{\pm.03}$ & 1.11$^{\pm.03}$
    & 132.2$^{\pm34.7}$ & 255.0$^{\pm.5}$ \\
  & ~~+no flex
    & 0.84$^{\pm.00}$
    & 231.9$^{\pm45.0}$ & 186.1$^{\pm.7}$
    & 1.75$^{\pm.03}$ & 1.23$^{\pm.03}$
    & 128.6$^{\pm33.5}$ & 275.5$^{\pm.8}$ \\
\bottomrule
\end{tabular}
\caption{Similarity distances: mean$^{\pm\text{95\% CI}}$ ($\downarrow$ better).
Best (lowest) value per (log, metric) pair underlined.
$^{\dagger}$DRL-DRL CTD on BPIC12\_W falls below DM-RR because
cycle-time compression narrows the distribution toward the
original's shape.}
\label{tab:similarity_drl}
\end{table*}

\noindent\textbf{Control-flow (NGD).}
DRL-DRL raises NGD to 0.46--0.75 relative to the DM-RR baseline of
0.18--0.36. The no-flex ablation isolates the mask's contribution:
removing it nearly doubles NGD on AC\_CRE (0.46\,$\to$\,0.89) and
raises it substantially on BPIC17\_W (0.68\,$\to$\,0.84) with a
CIR difference of at most 0.001, demonstrating that the mask halves
control-flow deviation at zero performance cost.
On BPIC12\_W the mask provides no benefit; NGD decreases without
it (0.75\,$\to$\,0.64), consistent with the Markov calibration
issue identified in Section~\ref{subsec:performance}: when the
routing prior encodes artificial loops, the mask anchors to a
distorted reference and the unmasked agent finds cleaner paths.

\noindent\textbf{Activity timing (AED, RED).}
Both metrics rise for DRL-DRL in proportion to the cycle-time
compression achieved, most clearly on BPIC17\_W where RED reaches
179.1 against a DM-RR baseline of 11.5. The no-flex ablation
produces near-identical values, confirming that AED and RED reflect
compression rather than mask behavior.

\noindent\textbf{Circadian, workforce, and congestion (CED, CWD, CAR).}
CED, CWD, and CAR remain stable across all policies and both DRL
variants, confirming that the agent optimizes within the temporal
and workforce constraints of the original process. The no-flex
ablation produces comparable values for all three metrics,
indicating that these dimensions are governed by the simulator's
temporal model rather than the activity mask.

\noindent\textbf{Cycle-time distribution (CTD).}
CTD rises for DRL-DRL on AC\_CRE and BPIC17\_W, the direct
signature of successful compression. On BPIC12\_W it falls below
DM-RR (146.5 vs.\ 170.3) because extreme compression narrows the
distribution toward the original log's shape.

% ---- 6.4  THREATS TO VALIDITY --------------------------------
\subsection{Threats to Validity}
\label{subsec:threats}

\noindent\textbf{Internal.}
The simulator supports XOR-only routing; processes with parallel
gateway structures are approximated as exclusive choices, which
distorts the branching probabilities the flexibility mask relies
on. This limitation is directly observable in the BPIC12\_W
results, where all Markov-based policies underperform random routing
and the mask provides no NGD benefit. Flexibility parameters were
selected empirically per log; a principled selection criterion
remains future work.

\noindent\textbf{External.}
Three logs from two domains are evaluated; applicability to
substantially different process structures is not guaranteed.

\noindent\textbf{Construct.}
The
\citeauthor{framework_for_measuring_the_quality_of_business_process_simulation_models}
similarity framework has been independently validated but may not
capture all dimensions of behavioral fidelity in every domain.
Stochasticity is mitigated by $K=10$ runs per (log, policy) with
95\%~CI reported throughout.

\section{Regularization Evaluation}
\label{sec:reg_evaluation}

This section evaluates RQ4: whether the policy-level regularization
term of Section~\ref{subsec:regularization_method} narrows the control-flow gap
the routing agent introduces while preserving its SLA gains. We reuse
the datasets, training configuration, and metrics of
Section~\ref{sec:flex_evaluation} without modification, and isolate the
term's effect through a focused ablation rather than a comparison
against additional policies.

\subsection{Experimental Setup}
\label{subsec:reg_setup}

The evaluation crosses two binary factors over the full joint agent
(DRL-DRL): the flexibility mask (on/off) and the KL regularization term
(on/off), yielding four configurations per log. Two of these, the
masked agent (\textit{flex}) and its unmasked ablation
(\textit{no flex}), are the unregularized results already reported in
Section~\ref{sec:flex_evaluation} and are carried over unchanged; the
remaining two, denoted \textit{+KL}, add the regularization term to each
mask setting. All other conditions are held fixed: the SLA training
threshold $T_{75}$, the flexibility parameters of
Table~\ref{tab:training}, and the $K=10$ simulated logs reported as
mean $\pm$ 95\%~CI.

The regularization strength $\lambda_{\text{KL}}$ is tuned per log and
held constant across both mask settings for that log
(Table~\ref{tab:reg_lambda}), so that the comparison between regularized
and unregularized variants is not confounded by a change in
$\lambda_{\text{KL}}$.

\begin{table}[!ht]
\centering\small
\setlength{\tabcolsep}{5pt}
\begin{tabular}{lcc}
\toprule
\textbf{Log} & \textbf{Threshold} & $\lambda_{\text{KL}}$ \\
\midrule
AC\_CRE   & $T_{75}$ & 0.08 \\
BPIC12\_W & $T_{75}$ & 0.01 \\
BPIC17\_W & $T_{75}$ & 0.05 \\
\bottomrule
\end{tabular}
\caption{Regularization strength per log.}
\label{tab:reg_lambda}
\end{table}

\subsection{Performance Analysis}
\label{subsec:reg_performance}

Table~\ref{tab:reg_compliance} reports CR and CIR at $T_{75}$ for the
four configurations.

\begin{table}[!ht]
\centering\scriptsize
\setlength{\tabcolsep}{2.5pt}
\renewcommand{\arraystretch}{1.1}
\begin{tabular}{ll cccc}
\toprule
& & \multicolumn{2}{c}{\textbf{No flex}} &
    \multicolumn{2}{c}{\textbf{Flex}} \\
\cmidrule(lr){3-4}\cmidrule(lr){5-6}
\textbf{Log} & \textbf{Met.} &
\textbf{base} & \textbf{+KL} &
\textbf{base} & \textbf{+KL} \\
\midrule
\multirow{2}{*}{AC\_CRE}
  & CR  & 99.5{\tiny(±0.2)} & 93.1{\tiny(±0.6)} & 99.4{\tiny(±0.2)} & 97.8{\tiny(±0.4)} \\
  & CIR & +32.9{\tiny(±0.3)} & +24.3{\tiny(±0.8)} & +32.8{\tiny(±0.3)} & +30.7{\tiny(±0.6)} \\
\midrule
\multirow{2}{*}{BPIC12\_W}
  & CR  & 100.0{\tiny(±0.0)} & 100.0{\tiny(±0.0)} & 99.9{\tiny(±0.0)} & 99.5{\tiny(±0.1)} \\
  & CIR & +33.3{\tiny(±0.0)} & +33.3{\tiny(±0.0)} & +33.2{\tiny(±0.1)} & +32.7{\tiny(±0.1)} \\
\midrule
\multirow{2}{*}{BPIC17\_W}
  & CR  & 100.0{\tiny(±0.0)} & 99.8{\tiny(±0.0)} & 100.0{\tiny(±0.0)} & 99.8{\tiny(±0.1)} \\
  & CIR & +33.3{\tiny(±0.0)} & +33.0{\tiny(±0.0)} & +33.3{\tiny(±0.0)} & +33.1{\tiny(±0.1)} \\
\bottomrule
\end{tabular}
\caption{CR (\%) and CIR (\%) at $T_{75}$ across the mask $\times$ KL
configurations of DRL-DRL. Values: mean (±95\% CI). Higher is better
($\uparrow$).}
\label{tab:reg_compliance}
\end{table}

\noindent\textbf{SLA compliance (RQ4).}
The regularization term preserves SLA compliance, and the cost it does
incur is consistently smaller when the mask is active. On the masked
agent, CIR at $T_{75}$ falls by at most $2.1$\,pp across all three logs
(AC\_CRE $+32.8 \to +30.7$, BPIC12\_W $+33.2 \to +32.7$, BPIC17\_W
$+33.3 \to +33.1$), with CR remaining near-saturated ($\geq 97.8\%$).
On the unmasked agent the same term costs more, most visibly on
AC\_CRE ($+32.9 \to +24.3$, a drop of $8.6$\,pp). This contrast is
expected: the mask already removes the implausible fast paths, so the
conformance pressure and the SLA objective compete less directly, and
the masked agent can satisfy the regularization term without
surrendering compliance. The mechanisms contribute on different axes rather than redundantly on the same one: the regularization term carries the control-flow reduction, while the mask, where its prior is faithful, confines the search to a trusted region and lowers the compliance cost of
regularization.

\subsection{Similarity Analysis}
\label{subsec:reg_similarity}

Table~\ref{tab:reg_similarity} reports the seven similarity metrics for
the four configurations.

\begin{table*}[!ht]
\centering\scriptsize
\setlength{\tabcolsep}{3pt}
\renewcommand{\arraystretch}{1.1}
\begin{tabular}{ll ccccccc}
\toprule
& & \multicolumn{7}{c}{\textbf{Metric}} \\
\cmidrule(lr){3-3}\cmidrule(lr){4-5}\cmidrule(lr){6-7}\cmidrule(lr){8-9}
\textbf{Log} & \textbf{Config}
  & \textbf{NGD}
  & \textbf{AED} & \textbf{RED}
  & \textbf{CED} & \textbf{CWD}
  & \textbf{CAR} & \textbf{CTD} \\
\midrule
\multirow{4}{*}{AC\_CRE}
  & no flex   & 0.89{\tiny(±.01)} & 65.6{\tiny(±11.1)} & 71.3{\tiny(±1.8)} & 2.33{\tiny(±.22)} & 2.35{\tiny(±.20)} & 105.6{\tiny(±18.1)} & 176.2{\tiny(±1.5)} \\
  & ~~+KL     & 0.34{\tiny(±.01)} & 65.5{\tiny(±9.5)} & 47.2{\tiny(±2.2)} & 2.48{\tiny(±.14)} & 2.50{\tiny(±.15)} & 92.1{\tiny(±16.5)} & 111.1{\tiny(±4.6)} \\
  & flex      & 0.46{\tiny(±.00)} & 70.1{\tiny(±16.2)} & 70.4{\tiny(±.7)} & 2.71{\tiny(±.19)} & 2.68{\tiny(±.20)} & 88.3{\tiny(±15.4)} & 165.8{\tiny(±1.6)} \\
  & ~~+KL     & \underline{0.32{\tiny(±.01)}} & 66.2{\tiny(±13.9)} & 60.4{\tiny(±1.6)} & 2.62{\tiny(±.18)} & 2.58{\tiny(±.19)} & 89.8{\tiny(±15.1)} & 142.0{\tiny(±3.2)} \\
\midrule
\multirow{4}{*}{BPIC12\_W}
  & no flex   & 0.64{\tiny(±.00)} & 252.0{\tiny(±23.4)} & 166.4{\tiny(±.4)} & 2.76{\tiny(±.10)} & 2.17{\tiny(±.11)} & 170.9{\tiny(±23.3)} & 145.9{\tiny(±.4)} \\
  & ~~+KL     & \underline{0.60{\tiny(±.00)}} & 249.4{\tiny(±24.5)} & 164.4{\tiny(±.5)} & 2.49{\tiny(±.06)} & 2.11{\tiny(±.10)} & 171.8{\tiny(±23.4)} & 143.9{\tiny(±.5)} \\
  & flex      & 0.75{\tiny(±.00)} & 253.6{\tiny(±23.8)} & 165.0{\tiny(±.3)} & 2.70{\tiny(±.08)} & 2.18{\tiny(±.13)} & 176.2{\tiny(±23.3)} & 146.5{\tiny(±.1)} \\
  & ~~+KL     & 0.69{\tiny(±.00)} & 244.8{\tiny(±24.2)} & 154.0{\tiny(±1.2)} & 2.60{\tiny(±.14)} & 2.08{\tiny(±.12)} & 175.7{\tiny(±23.3)} & 139.9{\tiny(±.7)} \\
\midrule
\multirow{4}{*}{BPIC17\_W}
  & no flex   & 0.84{\tiny(±.00)} & 231.9{\tiny(±45.0)} & 186.1{\tiny(±.7)} & 1.75{\tiny(±.03)} & 1.23{\tiny(±.03)} & 128.6{\tiny(±33.5)} & 275.5{\tiny(±.8)} \\
  & ~~+KL     & 0.52{\tiny(±.00)} & 217.8{\tiny(±44.7)} & 167.3{\tiny(±.4)} & 1.59{\tiny(±.03)} & 1.06{\tiny(±.02)} & 131.3{\tiny(±34.5)} & 242.8{\tiny(±.6)} \\
  & flex      & 0.68{\tiny(±.00)} & 229.8{\tiny(±44.1)} & 179.1{\tiny(±.2)} & 1.67{\tiny(±.03)} & 1.11{\tiny(±.03)} & 132.2{\tiny(±34.7)} & 255.0{\tiny(±.5)} \\
  & ~~+KL     & \underline{0.50{\tiny(±.00)}} & 216.1{\tiny(±45.2)} & 166.7{\tiny(±.5)} & 1.66{\tiny(±.03)} & 1.05{\tiny(±.03)} & 131.7{\tiny(±34.5)} & 238.5{\tiny(±.5)} \\
\bottomrule
\end{tabular}
\caption{Similarity distances across the mask $\times$ KL
configurations of DRL-DRL: mean$^{\pm\text{95\% CI}}$ ($\downarrow$
better). Best NGD per log underlined.}
\label{tab:reg_similarity}
\end{table*}

\noindent\textbf{Control-flow (NGD).}
The regularization term reduces control-flow deviation in every cell of
the grid, against both mask settings and on all three logs; its
benefit is therefore independent of the mask rather than contingent on
it. The reduction is largest where the routing prior the term targets
is most faithful: on AC\_CRE the masked agent falls from $0.46$ to
$0.32$ ($-30\%$) and the unmasked agent from $0.89$ to $0.34$
($-62\%$); on BPIC17\_W, from $0.68$ to $0.50$ ($-26\%$) and from
$0.84$ to $0.52$ ($-38\%$). On BPIC12\_W, where the XOR-only
approximation of parallel structures distorts the branching
probabilities, the reductions are smaller ($0.75 \to 0.69$ masked,
$0.64 \to 0.60$ unmasked), confirming that the term's effectiveness
tracks the fidelity of its reference. Notably, the term still delivers
a measurable gain on BPIC12\_W despite this structural limitation.

\noindent\textbf{Interaction with the mask.}
Both mechanisms read from the same empirical routing model, but they
apply it differently: the mask is a hard constraint that removes
activities outright, whereas the regularization term is a soft gradient
that the SLA objective can override where deviation is warranted. This
difference makes the mask's value conditional on the fidelity of the
routing model, while the term remains robust to it. On AC\_CRE, whose
prior is faithful, the hard constraint confines the search to a genuinely
good region and yields a strictly better policy: the masked, regularized
agent dominates its unmasked counterpart on both axes at once, reaching
NGD $0.32$ versus $0.34$ at CIR $+30.7$ versus $+24.3$, with the per-run
ranges essentially disjoint. On BPIC12\_W the ordering inverts, the best
control-flow similarity is reached without the mask
(\textit{no flex+KL}, $0.60$ versus $0.69$), because there the mask
enforces a distorted, XOR-approximated prior and removes routes the
regularizer would otherwise use, whereas the soft term merely bends the
policy toward the same flawed prior without being trapped by it. These
observations, across one faithful, one distorted, and one intermediate
log, are consistent with the hypothesis that hard masking amplifies
routing-model error while soft regularization absorbs it.

\noindent\textbf{Remaining dimensions (AED, RED, CED, CWD, CAR, CTD).}
The temporal, resource, and congestion metrics are undisturbed by the
regularization term: AED, CED, CWD, and CAR remain within their
confidence intervals of the corresponding unregularized agent on every
log, confirming that the term acts on the control-flow dimension it
targets and leaves the simulator-governed dimensions intact. Where the
term does move these metrics, it moves them toward the reference: RED
and CTD improve on AC\_CRE (RED $70.4 \to 60.4$, CTD $165.8 \to 142.0$
masked) and BPIC17\_W (RED $179.1 \to 166.7$, CTD $255.0 \to 238.5$
masked), a secondary benefit of routing the agent along more
representative paths.

\subsection{Threats to Validity}
\label{subsec:reg_threats}

\noindent\textbf{Internal.}
The regularization strength $\lambda_{\text{KL}}$ is selected per log.
Automated hyperparameter selection is a standing concern across the
field \cite{PrPm_research_directions} and is not addressed here; a principled or transferable choice of
$\lambda_{\text{KL}}$ is left to future work.

\noindent\textbf{Construct.}
The term draws the policy toward the empirical first-order routing
distribution, which is itself an approximation of the true process
behavior; on logs where that approximation is weak, as with the XOR
representation of parallel structures in BPIC12\_W, the reference the
term targets is correspondingly distorted.

\noindent\textbf{External.}
As in Section~\ref{sec:flex_evaluation}, three logs from two domains
are evaluated, and applicability to substantially different process
structures is not guaranteed. Stochasticity is mitigated by $K=10$ runs
per configuration with 95\%~CI reported throughout.


\section{Conclusion and Future Work}
\label{sec:conclusion}

This work introduced a unified decision framework in which a DRL
agent jointly selects the next activity and its executing resource
as a single action $(a,r)$ at each decision point in a business
process simulation. The framework rests on three contributions: the
formalization of \textbf{flexibility}, the dynamic delimitation of
the action space, tunable between process conformance and operational
performance; operationalized through the flexible masked policy and
its three interpretable parameters ($k$, $p$, $p_{\mathrm{end}}^{\min}$);
and a \textbf{policy-level regularization term} that draws the agent
toward the as-is routing distribution by acting directly on the
training objective rather than the reward signal.

The evaluation on three real-life event logs provides evidence across
all four research questions. On \textbf{RQ1}, DRL-DRL achieves
near-perfect SLA compliance at $T_{75}$ (CIR\,+32.8\% to +33.3\%),
with strong performance at every other threshold due to intermediate
reward shaping that teaches continuous cycle-time reduction. On
\textbf{RQ3}, the step from resource-only to full joint learning
raises the CIR in every case, by +12.4~pp on AC\_CRE and +21.6~pp
on BPIC17\_W; the BPIC12\_W results, where every Markov-based
policy falls below the random baseline yet DRL-DRL recovers to
+33.2\%, demonstrate that joint optimization is structurally
enabling rather than merely additive.
On \textbf{RQ2}, the no-flex ablation shows that the mask halves NGD
deviation on logs with predominantly sequential structure at zero
compliance cost, while CED, CWD, and CAR remain stable across all
variants. On \textbf{RQ4}, the regularization term reduces
control-flow deviation in every configuration of the
mask\,$\times$\,KL grid and on every log, lowering NGD by up to 62\%
over the unregularized agent while preserving SLA compliance.

The two control-flow mechanisms relate to the routing model in
different ways. The mask is a hard constraint that removes activities
outright, whereas the regularization term is a soft gradient the SLA
objective can override where deviation is warranted. The mask's value
is therefore conditional on the fidelity of the routing model, while
the term is comparatively robust to it. Where the prior is faithful,
as on AC\_CRE, the hard constraint confines the search to a trusted
region and yields a strictly better policy, the masked, regularized
agent dominates its unmasked counterpart on both NGD and compliance at
once. Where the prior is distorted, as on BPIC12\_W under the XOR
approximation of parallel structures, the mask enforces a flawed region
and the regularization term alone performs best. Across one faithful,
one distorted, and one intermediate log, these observations are
consistent with the broader pattern that the effectiveness of every
conformance mechanism is bounded by the fidelity of the routing
representation it relies on.


Beyond its technical contributions, the framework yields two
complementary outputs that are worth distinguishing. The simulation
environment produces an optimized process log that can be used for
post-hoc analysis, scenario comparison, and process redesign. The
learned policy, however, is the more consequential artifact: it is
a portable, optimized decision function that, given the current state
of any ongoing case and the available resources, immediately
recommends the joint activity--resource action most likely to meet
the SLA. This is the sense in which the framework is most naturally
understood as a prescriptive process monitoring tool. Unlike
predictive approaches, which inform a manager that a case is at
risk, the policy tells them what to do about it. It functions as a
consultant that has internalized the dynamics of the process and can
be queried at any decision point during live execution, without
requiring a new simulation run. This positioning opens practical
deployment scenarios beyond the evaluation setting, including
real-time decision support in workflow management systems and
automated dispatching in high-volume operational environments.

Several directions extend the present work. First, extending the
simulator to parallel branching would remove the calibration
distortion observed on BPIC12\_W and is expected to improve both
performance and behavioral fidelity simultaneously. Second, both the
flexibility parameters and the regularization strength
$\lambda_{\text{KL}}$ are currently selected per log; a principled or
transferable selection criterion, and, more broadly, automated
hyperparameter tuning, would reduce the manual effort the framework
requires and is a natural next step. Third, replacing the engineered
state representation with a learned process embedding, for instance
via an LSTM or transformer encoder trained on case prefixes, could
allow the agent to capture latent process patterns and make more
informed decisions beyond what hand-crafted features provide. Closely
related, enriching the agent's context with information that cannot
be derived from the log alone, such as domain knowledge, resource
profiles, or process-level constraints encoded through an NLP model
and concatenated to the state, could allow the agent to reason over
richer operational contexts and produce more informed recommendations.
Fourth, an interesting dimension not yet addressed is the cost of
resources; the present framework treats all resources as equally
expensive, yet in practice their costs are heterogeneous. Extending
the approach through additional multi-objective regularization terms,
which have already demonstrated good results in the conformance
setting, offers a natural path toward cost-aware policies, and CWD
provides a suitable proxy metric for their evaluation. Fifth, the
admissible action space could be further constrained by
domain-specific operational rules, such as resource unavailability
or regulatory restrictions, enabling practitioners to query the agent
under ad-hoc what-if conditions and bringing the framework closer to
a practical decision-support tool. Finally, extending the evaluation
beyond distributional similarity toward an explainability analysis of
the agent's intervention patterns would assess whether the recommended
decisions are interpretable and actionable for human operators,
bridging prescriptive process monitoring and human-in-the-loop
process management.

\bigskip
\noindent\textit{Reproducibility.} The source code is available at
\url{https://github.com/AdaptiveBProcess/Optimal-Predictor-of-Resources-and-Activities}
and the event logs and aggregated results are archived at
\url{https://zenodo.org/records/20477596}.

\section*{References}
\label{sec:references}

\bibliographystyle{IEEEtranN}
\begingroup
\renewcommand{\section}[2]{}%
%\renewcommand{\chapter}[2]{}% for other classes
\small
\bibliography{main}
\endgroup


\end{document}

    